\documentclass[11pt,reqno]{amsart}
\usepackage[margin=1.05in]{geometry}
\usepackage{amsmath,amssymb,amsthm,mathtools,booktabs,longtable,array,enumitem}
\usepackage[hyphens]{url}
\usepackage[colorlinks=true,linkcolor=blue,citecolor=blue,urlcolor=blue]{hyperref}
\usepackage[expansion=false]{microtype}

\setlist{leftmargin=2em}
\emergencystretch=2em

\theoremstyle{plain}
\newtheorem{theorem}{Theorem}[section]
\newtheorem{proposition}[theorem]{Proposition}
\newtheorem{lemma}[theorem]{Lemma}
\newtheorem{corollary}[theorem]{Corollary}
\newtheorem{fact}[theorem]{Fact}
\newtheorem*{theoremA}{Theorem A}
\newtheorem*{theoremB}{Theorem B}
\theoremstyle{definition}
\newtheorem{conjecture}[theorem]{Conjecture}
\newtheorem{definition}[theorem]{Definition}
\newtheorem{remark}[theorem]{Remark}
\newtheorem{certificate}[theorem]{Certificate}

% Certificate line attached to a lemma proved by computation
\newcommand{\cert}[2]{\par\smallskip{\raggedright\noindent\emph{Certificate:} \texttt{#1}; log \texttt{#2}.\par}}

\DeclareMathOperator{\Li}{Li}
\DeclareMathOperator{\TV}{TV}
\DeclareMathOperator{\erfc}{erfc}
\DeclareMathOperator{\sgn}{sgn}
\DeclareMathOperator{\Amp}{Amp}
\renewcommand{\Re}{\operatorname{Re}}
\renewcommand{\Im}{\operatorname{Im}}
\newcommand{\Z}{\mathbb Z}
\newcommand{\N}{\mathbb N}
\newcommand{\C}{\mathbb C}
\newcommand{\R}{\mathbb R}
\newcommand{\HH}{\mathbb H}
\newcommand{\cP}{\mathcal P}
\newcommand{\cG}{\mathcal G}
\newcommand{\cL}{\mathcal L}
\newcommand{\cN}{\mathcal N}
\newcommand{\cB}{\mathcal B}
\newcommand{\fa}{\mathfrak a}
\newcommand{\Rh}{\widehat R}
\newcommand{\St}{\widetilde S}
\newcommand{\Qt}{\widetilde Q}
\newcommand{\Econ}{E_{\mathrm{con}}}
\newcommand{\e}{\mathrm e}
\newcommand{\vp}{\varpi}

\begin{document}

\title[A method for Borwein-type sign problems]{Sign patterns of Borwein products by explicit circle-method and saddle-point certificates: the cubic Borwein conjecture and the Third Borwein conjecture}
\author{Jaideep Sai Padhi}
\address{Department of Computer Science, Purdue University, West Lafayette, IN 47907, USA}
\email{jpadhi@purdue.edu}
\date{\today}
\subjclass[2020]{Primary 11P82; Secondary 05A17, 30E15, 65G20}
\keywords{Borwein conjecture, sign patterns, circle method, saddle-point method, ball arithmetic, computer-assisted proof}

\begin{abstract}
For a prime $p$ put $P_{p,n}(q)=\prod_{k\le pn,\,p\nmid k}(1-q^k)$. We develop a method for proving sign patterns of the
coefficients of $P_{p,n}(q)^\delta$ modulo $p$, and use it to settle two open cases. First, we prove the remaining residue
class of the cubic Borwein conjecture: $[q^m]P_{3,n}(q)^3\le0$ for all $n\ge1$ and all $m\equiv2\pmod 3$; together with
the work of Wang and Krattenthaler this proves the conjecture. Second, we prove the Third Borwein conjecture: the
coefficients of $P_{5,n}(q)=(q;q)_{5n}/(q^5;q^5)_n$ are $\ge0$ at exponents divisible by $5$ and $\le0$ otherwise.

The method factors $P_{p,n}^\delta=G_p^\delta E_n$ with $G_p=(q;q)_\infty/(q^p;q^p)_\infty$ and a tail $E_n$ with
non-negative coefficients, uses an explicit circle-method expansion of $G_p^\delta$ with a fully proved error constant, and
splits the coefficient range into an exact range, a band, a Laplace (saddle-point) region and a centre. Each analytic
region is reduced to finitely many inequalities in a few real parameters, which are verified in ball arithmetic uniformly
in $n$. For $p=5$ new difficulties arise: two residue classes in which the main term vanishes, which require re-centring
the saddle point on the thin factor, and a centre region that requires an off-peak bound for $|P_{5,n}|$ on the whole
circle.

\emph{Both proofs are computer-assisted.} All programs, their logs, and the lines each log must contain are published
with the paper and are described in \S\ref{sec:computations}.
\end{abstract}

\maketitle
\setcounter{tocdepth}{1}
\tableofcontents

%=====================================================================
\section{Introduction}\label{sec:intro}
%=====================================================================

For a prime $p$ and $n\ge1$ put
\[
P_{p,n}(q)=\prod_{\substack{k\le pn\\ p\nmid k}}(1-q^k)=\frac{(q;q)_{pn}}{(q^p;q^p)_n}.
\]
P.~Borwein observed, and Andrews \cite{Andrews} made widely known, that the coefficients of such products and of their
small powers appear to have a sign pattern that depends only on the exponent modulo $p$.
% Andrews [J. Symbolic Comput. 20 (1995)] discusses several conjectures of P. Borwein for p=3 and p=5; the names
% First/Second/Third follow Wang--Krattenthaler and Berkovich--Dhar (First: P_{3,n}; Second: P_{3,n}^2; Third: P_{5,n}).
% TODO: numbering inside Andrews's paper itself not verified (full text not accessible).
The three conjectures that concern us are the following. Write $c(m)$ for the coefficient of $q^m$ in the polynomial
under discussion.
\begin{enumerate}
\item (\emph{First Borwein conjecture.}) For $P_{3,n}$: $c(m)\ge0$ if $m\equiv0$ and $c(m)\le0$ if $m\equiv1,2\pmod 3$.
\item (\emph{Cubic Borwein conjecture.}) The same pattern $+\,-\,-$ holds for $P_{3,n}^3$.
\item (\emph{Third Borwein conjecture.}) For $P_{5,n}$: $c(m)\ge0$ if $5\mid m$ and $c(m)\le0$ otherwise.
\end{enumerate}
There is also a second Borwein conjecture, for $P_{3,n}^2$. (The names First, Second and Third follow
\cite{WK,BerkovichDhar}; Andrews \cite{Andrews} discusses Borwein's conjectures for $p=3$ and $p=5$.) The first conjecture was
proved by Wang \cite{Wang}, by an analytic (saddle-point) method. Wang and Krattenthaler \cite{WK} gave a new proof of the first
conjecture and proved the second by an asymptotic approach to Borwein-type sign patterns; they also proved ``two thirds'' of
the cubic conjecture (see also \cite[\S1]{BerkovichDhar}), namely the class $m\equiv0\pmod3$ and, in the range
$m\le\frac92n^2$ (equivalently, by the palindromic symmetry of $P_{3,n}^3$, the complementary halves), the class
$m\equiv1\pmod3$, and explain \cite[\S11]{WK} why the class $m\equiv2$ resists their method.
% Consistency check (resolves the earlier TODO): P_{3,n}^3 is palindromic of degree 9n^2 = 0 (mod 3), so m -> 9n^2-m
% swaps the classes 1 and 2. Berkovich--Dhar's 'all of kappa_n, half of delta_n, half of theta_n' therefore describes
% class 0 plus one pair of mirror-image halves (class 1 for m <= 9n^2/2 <-> class 2 for m >= 9n^2/2). The half that is
% open must be the 'thin' one: by Jacobi's identity every coefficient of P_infty^3 in class 2 vanishes, so class 2 with
% m <= 9n^2/2 is exactly where the main term cancels -- which is what WK Sect. 11 says resists their method. This is the
% class treated in Theorem A. Theorem number in WK Sect. 10 not checked verbatim (arXiv full text not reachable here).
The Third Borwein conjecture is listed as open in the recent literature \cite[Conj.~1.3]{BerkovichDhar}, where it is noted
that \cite{WK} provides ideas for ``three fifths'' of a proof; we are not aware of a proof in the literature.

\subsection{Main results}

\begin{conjecture}[cubic Borwein conjecture]\label{conj:cubic}
Let $c_n(m)=[q^m]P_{3,n}(q)^3$. For all $n\ge1$ and $0\le m\le\frac92n^2$: $c_n(m)\ge0$ if $m\equiv0$, and $c_n(m)\le0$ if
$m\equiv1,2\pmod 3$.
\end{conjecture}

\begin{theoremA}[the remaining class of the cubic conjecture]
For all $n\ge1$ and all $m\equiv2\pmod3$ with $m\le\frac92n^2$ we have $c_n(m)=0$ if $m<3n$ and $c_n(m)<0$ if $m\ge3n$.
Together with \cite{WK} this proves Conjecture~\ref{conj:cubic}.
\end{theoremA}

\begin{conjecture}[Third Borwein conjecture]\label{conj:third}
Let $c_n(m)=[q^m]P_{5,n}(q)$. For every $n\ge1$ and every $m$, $c_n(m)\ge0$ if $m\equiv0\pmod5$ and $c_n(m)\le0$ if
$m\not\equiv0\pmod 5$.
\end{conjecture}

\begin{theoremB}
Conjecture~\ref{conj:third} holds.
\end{theoremB}

Theorem~A is Theorem~\ref{thm:cubic-main} and Theorem~B is Theorem~\ref{thm:third-main}. Both are computer-assisted.
Theorem~A needs only the one residue class $m\equiv2\pmod3$; Theorem~B is proved for all five classes modulo $5$.

\subsection{Why the remaining cases are hard}
Both polynomials factor as $P_{p,n}^\delta=G_p^\delta E_n$ (\S\ref{sec:framework}), where $G_p=(q;q)_\infty/(q^p;q^p)_\infty$
and $E_n$ is a tail product with non-negative coefficients. The circle method gives the coefficients of $G_p^\delta$ with a
main term from the cusps of denominator $p$, and the tail $E_n$ is a small perturbation of $1$ for exponents $m$ below
$N=pn+1$. The difficulty is that in the classes we must treat, the main term of $G_p^\delta$ vanishes:
\begin{itemize}
\item for $p=3$, $\delta=3$, all coefficients of $G_3^3$ in the class $m\equiv2$ vanish (Lemma~\ref{lem:vanish});
\item for $p=5$, $\delta=1$, all coefficients of $G_5$ in the classes $m\equiv3,4$ vanish (Lemma~\ref{lem:vanish}).
\end{itemize}
In these ``thin'' classes $c_n(m)$ is produced entirely by the tail $E_n$, and the leading saddle-point contribution
cancels. The sign then comes from a secondary term that is exponentially smaller than the naive main term.

\subsection{Overview of the method}
The coefficient range is split into four regions (\S\ref{sec:regions}).
\begin{enumerate}
\item \emph{Exact range.} Small $n$ are handled by exact integer computation.
\item \emph{Band.} For $m$ below a small multiple of $N$ only a few parts of $E_n$ contribute, and $c_n(m)$ is a short,
explicit combination of coefficients of $G_p^\delta$, controlled by exact values and by the circle-method expansion.
\item \emph{Laplace region.} Here an exact Bessel--Laplace (or Poisson) representation of the main term is combined with
the closed form of $\log E_n$ at a twisted point, and a saddle-point analysis is carried out in scaled variables.
\item \emph{Centre.} Near the middle of the coefficient range, Cauchy's formula for the polynomial is evaluated directly
by steepest descent through the saddles near the primitive $p$-th roots of unity.
\end{enumerate}
The method was developed on the cubic case. For $p=5$ three new difficulties arise:
\begin{itemize}
\item \emph{two} vanishing classes instead of one;
\item in the Laplace region for the classes $3,4$ the surviving term carries a factor $e^{-5nu}$ which is of size
$e^{O(N)}$ \emph{in the exponent}; the saddle must be re-centred on this ``thin factor'' (\S\ref{sec:lap34});
\item in the centre, the contribution of the circle away from the four peaks must be bounded by an explicit
\emph{off-peak lemma} for $|P_{5,n}|$ (\S\ref{sec:offpeak}), which is then combined with the near-peak analysis box by box.
\end{itemize}
The two proofs share the method, not their results: the cubic theorem is not an ingredient of the proof for $p=5$, and
conversely.

\subsection{What is computer-assisted, and how it is verified}
Each analytic region is reduced to finitely many inequalities between explicit functions of one to three real parameters
(a scaled saddle point, the reciprocal $1/N$, an angle), and these are verified on coverings of the parameter domains by
balls in Arb ball arithmetic \cite{Arb}. Uniformity in $n$ is obtained in two ways: the scaled integrands are analytic in
$1/N$ at $0$, so one ball covers all large $N$; and every term that depends on $n$ is shown, by an asserted inequality, to
be non-increasing in $n$. Exact ranges are computed with FLINT \cite{FLINT} in $\Z[q]$.

Throughout, a statement proved by a computation is stated as a lemma or a certificate, followed by a line
\emph{Certificate:} naming the program and its log. Section~\ref{sec:computations} lists, for every certificate, the
command, the log and the line the log must contain; Appendix~\ref{app:verification} summarises the independent audits of the
proof of Theorem~B. The trust base consists of the written arguments of this paper, the standard facts of
Appendix~\ref{app:facts}, and the correctness of Arb and FLINT.

\subsection{Organisation}
Section~\ref{sec:framework} sets up the common framework and notation. Section~\ref{sec:cubic} proves Theorem~A and
Section~\ref{sec:third} proves Theorem~B. Section~\ref{sec:computations} describes the certificates.

%=====================================================================
\section{The general framework}\label{sec:framework}
%=====================================================================

Throughout, $(p,\delta)$ is one of the two cases
\[
(p,\delta)=(3,3)\quad\text{(Section~\ref{sec:cubic})},\qquad (p,\delta)=(5,1)\quad\text{(Section~\ref{sec:third})},
\]
and $c_n(m)=[q^m]P_{p,n}(q)^\delta$.

\subsection{Notation}\label{sec:notation}
Table~\ref{tab:notation} fixes the symbols shared by the two parts. Where the two source proofs used one letter for
different objects, one of them has been renamed; symbols that are local to a single subsection are defined where they are
used and are not listed.

\begin{table}[h]
\centering
\small
\renewcommand{\arraystretch}{1.15}
\begin{tabular}{@{}p{0.23\textwidth}p{0.72\textwidth}@{}}
\toprule
symbol & meaning\\
\midrule
$P_{p,n}$; $P_n$ & $\prod_{k\le pn,\,p\nmid k}(1-q^k)$; in \S\ref{sec:cubic} and \S\ref{sec:third} $p$ is fixed and we write $P_n$\\
$c_n(m)$ & $[q^m]P_{p,n}^\delta$; $\deg P_{p,n}^\delta=\delta(p-1)pn^2/2$ ($9n^2$ for $p=3$, $10n^2$ for $p=5$)\\
$G_p$; $s(i)$ & $(q;q)_\infty/(q^p;q^p)_\infty=\prod_{p\nmid k}(1-q^k)$; $s(i)=[q^i]G_p^\delta$\\
$E_n$; $e(j)$ & $\prod_{k>pn,\,p\nmid k}(1-q^k)^{-\delta}$, so that $P_{p,n}^\delta=G_p^\delta E_n$; $e(j)=[q^j]E_n\ge0$\\
$N$ & $pn+1$ (the smallest part of $E_n$); in \S\ref{sec:G5} only, $N$ is the Farey order\\
$\omega_p$ & $e^{2\pi i/p}$ (so $\omega_3=\omega$ of \S\ref{sec:cubic}, $\omega_5$ of \S\ref{sec:third})\\
$B$, $C$ & $B=\delta(p-1)/12$ ($=\frac12$ resp.\ $\frac13$), $C=C(i)=2i-B$\\
$\cP_k(i)$ & $\frac{2\pi}{k}\sqrt{B/C}\,I_1\big(\frac{2\pi}{k}\sqrt{BC}\big)$ for $i\ge1$; $\cP_k(i)=0$ for $i\le0$\\
$a_k$; $a_p$ & $a_k=2\pi^2B/k^2$; so $a_3=\pi^2/9$ (cubic), $a_5=2\pi^2/75$ (Third), and
$\cP_p(i)=\sqrt{a_p/x}\,I_1(2\sqrt{a_px})$ with $x=i-B/2$\\
$\fa(i)$ & the amplitude of the $k=p$ main term: $\fa(i)=\sqrt3\chi(i)$ for $p=3$; $\fa(i)=2\cos(2\pi i/5+\pi/5)+2\cos(4\pi i/5)$ for $p=5$\\
$K(i)$ & $\lfloor\sqrt{4\pi i}\rfloor+2$\\
$\Econ$ & error constant of the expansion: $50.1501$ ($p=3$), $27.4636$ ($p=5$)\\
$\mu$ & $m/\deg P_{p,n}^\delta$, so the half-range is $\mu\le\frac12$ (used in \S\ref{sec:third});
in \S\ref{sec:cubic} the centre is parametrised by $\bar\mu=m/N^2$ and $\tilde\mu=(m-N-\frac14)/N^2$\\
$W$ & real part of the Laplace variable ($w=W+iy$ in \S\ref{sec:cubic}, $u=W(1+i\eta)$ in \S\ref{sec:third})\\
$v$; $V$ & scaled saddle: $v=NW$ in \S\ref{sec:cubic}; $V=pnW=5nW$ in \S\ref{sec:third}\\
$\lambda$ & $m/N$ (\S\ref{sec:third}; the band is $\lambda<4$)\\
$\cL$; $L$ & $\cL(u)=3\sum_{k\le3n,3\nmid k}\log(1-e^{ku})$ (cubic centre); $L$ the radial parameter $r=e^{-L/(5n)}$
(Third centre)\\
$\Li_s$ & polylogarithm, $\Li_1(x)=-\log(1-x)$\\
$B_1,B_2$; $\tilde B_k$ & $B_1(x)=x-\frac12$, $B_2(x)=x^2-x+\frac16$; $\tilde B_k$ their $1$-periodic extensions\\
class & $m\bmod p$\\
\bottomrule
\end{tabular}
\caption{Notation shared by Sections~\ref{sec:cubic} and~\ref{sec:third}.}
\label{tab:notation}
\end{table}

\subsection{Factorisation, palindromy and vanishing classes}

\begin{lemma}[factorisation]\label{lem:factor}
$P_{p,n}^\delta=G_p^\delta\,E_n$. The series $E_n$ has non-negative coefficients, $e(0)=1$ and $e(j)=0$ for $0<j<N$.
Consequently
\[
c_n(m)=\sum_{j\ge0}e(j)\,s(m-j)=s(m)+\sum_{j\ge N}e(j)s(m-j).
\]
\end{lemma}
\begin{proof}
$G_p=\prod_{p\nmid k}(1-q^k)$. Split this product at $k=pn$: the factor over $k>pn$ is $E_n^{-1/\delta}$, and $E_n$ is a
product of ($\delta$-th powers of) geometric series in $q^k$ with $k\ge pn+1=N$.
\end{proof}

\begin{lemma}[palindromy]\label{lem:palin}
$P_{p,n}^\delta$ is palindromic of degree $D=\delta(p-1)pn^2/2$. The map $m\mapsto D-m$ preserves the class $0$ and maps
the class $s$ to $-s$. Hence the sign patterns of Conjectures~\ref{conj:cubic} and~\ref{conj:third} need only be proved
for $0\le m\le D/2$.
\end{lemma}
\begin{proof}
$P_{p,n}$ is a product of $(p-1)n$ factors $1-q^k$, each satisfying $q^k(1-q^{-k})=-(1-q^k)$, and
$\sum_{k\le pn,\,p\nmid k}k=\frac{pn(pn+1)}2-\frac{p\,n(n+1)}2=\frac{(p-1)pn^2}{2}$. For $p=3$ the sign is
$(-1)^{2n\cdot3}=1$ and $D=9n^2$; for $p=5$ it is $(-1)^{4n}=1$ and $D=10n^2$. In both cases $p\mid D$. For $p=3$ the
classes $1,2$ are exchanged and carry the same conjectured sign; for $p=5$ the classes $1\leftrightarrow4$ and
$2\leftrightarrow3$ are exchanged, and again the conjectured sign is invariant under $s\mapsto-s$.
\end{proof}

\begin{lemma}[vanishing classes]\label{lem:vanish}
\begin{enumerate}
\item[(a)] ($p=3$, $\delta=3$.) $s(i)=[q^i]G_3^3=0$ for $i\equiv2\pmod3$.
\item[(b)] ($p=5$, $\delta=1$.) $s(i)=[q^i]G_5=0$ for $i\equiv3,4\pmod5$.
\end{enumerate}
\end{lemma}
\begin{proof}
(a) By Jacobi's identity $(q;q)_\infty^3=\sum_{j\ge0}(-1)^j(2j+1)q^{j(j+1)/2}$, and $j(j+1)/2\equiv0,1,0\pmod3$ for
$j\equiv0,1,2$. So $(q;q)_\infty^3$ is supported on the classes $\{0,1\}$, and $G_3^3=(q;q)^3_\infty/(q^3;q^3)^3_\infty$ is
this series times a series in $q^3$.

(b) By Euler's pentagonal theorem, $(q;q)_\infty=\sum_{k\in\Z}(-1)^kq^{k(3k-1)/2}$. Modulo $5$ the exponent $k(3k-1)/2$
takes the values $0,1,0,2,2$ for $k\equiv0,1,2,3,4$. So $(q;q)_\infty$ is supported on the residues $\{0,1,2\}$, and
multiplying by $1/(q^5;q^5)_\infty$, a series in $q^5$, preserves this.
\end{proof}

In the vanishing (``thin'') classes the main term of $s(m)$ is absent and the sign of $c_n(m)$ is produced by the
convolution with $E_n$.

\subsection{The circle-method expansion of the eta-quotient}\label{sec:circle-general}

Let $f(x)=\prod_{m\ge1}(1-x^m)^{-1}$, $\tau=h/k+iz/k^2$ with $\Re z>0$, $(h,k)=1$, and $d=\gcd(p,k)$. The eta
transformation, in the form of \cite[Prop.~7]{SZ} with the Hardy--Ramanujan convention $hh'\equiv-1\pmod k$
\cite{Padhi_SZ}, gives
\begin{equation}\label{eq:Gp-transform}
G_p(e^{2\pi i\tau})^\delta=\Big(\frac pd\Big)^{\delta/2}e^{\pi i\delta(-s(h,k)+s(ph/d,k/d))}
\exp\Big(\frac{\pi\delta(d^2-p)}{12pz}-\frac{(p-1)\pi\delta z}{12k^2}\Big)\widehat G^\delta,
\end{equation}
where
\[
\widehat G=\frac{f(e^{2\pi idh'_d/k-2\pi d^2/(pz)})}{f(e^{2\pi ih'/k-2\pi/z})},
\]
with $s(h,k)$ the Dedekind sum and principal powers. For our two cases this has the following structure.
\begin{itemize}
\item \emph{Growing cusps} ($p\mid k$, $d=p$). The prefactor is unimodular and the exponential is
$\exp(\pi B/z-\pi Bz/k^2)$ with $B=\delta(p-1)/12$. The first correction term of $\widehat G^\delta$ grows only if
$\delta>24/(p-1)$, which is not the case for $(3,3)$ or $(5,1)$. Integrated over the Ford circle, the leading term gives
the Bessel main term $\cP_k(i)$, since $\frac{i}{k^2}\oint_K\exp(\pi B/z+\pi Cz/k^2)\,dz=\cP_k(i)$ for the circle
$K=\{|z-\frac12|=\frac12\}$ oriented clockwise (Lemma~\ref{lem:bessel-contour}).
\item \emph{Bounded cusps} ($p\nmid k$). On $\Re(1/z)\ge1$, $|G_p^\delta|\le K_p^\delta$ with
$K_p=\sqrt p\,e^{-\pi(p-1)/(12p)}f(e^{-2\pi/p})f(e^{-2\pi})$.
\end{itemize}
Carrying out the Rademacher dissection with Farey order at least $\sqrt{4\pi i}+1$ and moving each integral to the chord
of its Ford arc gives an expansion of the form
\begin{equation}\label{eq:expansion-shape}
s(i)=\fa(i)\,\cP_p(i)+\rho(i),\qquad
|\rho(i)|\le\sum_{\substack{2p\le k\le K(i)\\ p\mid k}}\varphi(k)\cP_k(i)+\Econ ,
\end{equation}
with an explicit constant $\Econ=E_{\rm arc}+E_{\rm rem}+E_{\rm ng}$ (arc completion, remainder at growing cusps, bounded
cusps). The two instances are Theorem~\ref{thm:expansion3} ($p=3$; proved in \S\ref{sec:expansion3}) and
Theorem~\ref{thm:G5} ($p=5$; proved in full in \S\ref{sec:G5}). The main-term amplitude $\fa(i)$ at the cusps of
denominator $p$ is computed exactly from the Dedekind sums $s(h,p)$, and it vanishes precisely on the thin classes of
Lemma~\ref{lem:vanish}.

\subsection{The region decomposition}\label{sec:regions}

For each $n$ the half-range $0\le m\le D/2$ is covered by the following regions (the precise thresholds are in
Tables~\ref{tab:regions3} and~\ref{tab:regions5}).
\begin{enumerate}
\item \textbf{Exact range}, $n\le n_0$: exact computation of $P_{p,n}^\delta$ in $\Z[q]$.
\item \textbf{Band}, $m<cN$ with $c=2$ ($p=3$) or $c=4$ ($p=5$): since $e(j)=0$ for $0<j<N$, only partitions of $j$
into at most $c-1$ parts from $E_n$ contribute, and $c_n(m)$ is a short explicit sum of values $s(i)$.
\item \textbf{Laplace region}: the main term of~\eqref{eq:expansion-shape} is represented as a Laplace-type integral
(Poisson form for $p=3$, Bessel--Laplace form for $p=5$), convolved with $E_n$ at a twisted point, and evaluated by a
saddle-point analysis in scaled variables. The remainder $\rho$ is bounded by Rankin's trick at the same saddle.
\item \textbf{Centre}: Cauchy's formula for the polynomial, evaluated through the saddles near the primitive $p$-th roots
of unity, together with bounds for the rest of the circle.
\end{enumerate}

\subsection{Ball-arithmetic certificates}\label{sec:certs-general}

Every inequality on which an analytic region rests is reduced to finitely many inequalities between explicit functions
of a few real parameters and verified in ball arithmetic (Arb \cite{Arb}, via \texttt{python-flint}). A ball computation
returns an enclosure of the exact value over the whole parameter ball; an inequality is \emph{certified} on a ball when
the enclosure lies strictly on the correct side. The parameter domains are covered by finitely many balls (tiles);
failing balls are bisected. Every program aborts on a failed assertion and prints a summary line, and it is that line
which constitutes the certificate (\S\ref{sec:computations}).

Uniformity in $n$ uses two devices.
\begin{itemize}
\item \emph{Analyticity in $1/N$.} After scaling, the integrands are analytic in $\epsilon=1/N$ at $\epsilon=0$, so one
ball $\epsilon\in[0,\epsilon_1]$ covers all $N\ge1/\epsilon_1$ (\S\ref{sec:cubic}).
\item \emph{Monotonicity.} Every term depending on $n$ (or on a scaled size parameter $X$) is of the form
$n^ae^{-nb}$, or a finite sum of such terms, and a certified inequality on the exponents shows that it is non-increasing
in $n$ beyond the base point. One evaluation at the base point then covers all larger $n$.
\end{itemize}

%=====================================================================
\section{The cubic case: the class \texorpdfstring{$m\equiv2\pmod 3$}{m = 2 mod 3}}\label{sec:cubic}
%=====================================================================

In this section $p=\delta=3$, $P_n=P_{3,n}$, $c_n(m)=[q^m]P_n^3$, $N=3n+1$, $a_3=\pi^2/9$, $\omega=\omega_3=e^{2\pi i/3}$,
$s(i)=[q^i]G_3^3$ and $P_\infty=G_3$, so that $E_n=P_\infty^3/P_n^3=\prod_{k>3n,\,3\nmid k}(1-q^k)^{-3}$. By
Lemma~\ref{lem:vanish}(a), $c_n(m)$ for $m\equiv2$ is produced entirely by the tail $E_n$.

\begin{theorem}\label{thm:cubic-main}
For all $n\ge1$ and all $m\equiv2\pmod3$ with $m\le\frac92n^2$ we have $c_n(m)=0$ if $m<3n$ and $c_n(m)<0$ if $m\ge3n$.
Together with \cite{WK} this proves Conjecture~\ref{conj:cubic}.
\end{theorem}

Put $\bar\mu=m/N^2$ and $\tilde\mu=(m-N-\frac14)/N^2$. The admissible pairs $(n,m)$ are covered by five regions
(Table~\ref{tab:regions3}).

\begin{table}[h]
\centering
\small
\begin{tabular}{lll}
\toprule
region & method & where\\
\midrule
$n\le243$ & exact integer computation & \S\ref{sec:c-finite}\\
$n\ge244$, $m<6n+2$ & band lemma & Prop.~\ref{prop:c-band}\\
$n\ge244$, $m\ge 6n+2$, $\tilde\mu<0.0985$, saddle $v\le3000$ & Laplace certificate, uniform in $n$ & \S\ref{sec:c-laplace}\\
$n\ge244$, saddle $v\ge3000$ & closing lemma & \S\ref{sec:c-closing}\\
$n\ge244$, $\bar\mu\ge0.098$ & centre certificate, uniform in $n$ & \S\ref{sec:c-centre}\\
\bottomrule
\end{tabular}
\caption{The covering of the admissible $(n,m)$ for $p=3$. That these regions exhaust all cases is
Proposition~\ref{prop:c-coverage}.}
\label{tab:regions3}
\end{table}

\subsection{Structural lemmas}\label{sec:c-structural}

\begin{lemma}[tail phase]\label{lem:c-tailphase}
For $0<q<1$, $\Delta(q):=\Im\log E_n(\omega q)\in(0,\pi/2)$.
\end{lemma}
\begin{proof}
$\Im\log E_n(\omega q)=3\sum_{A\ge n}\big[f(q^{3A+1})-f(q^{3A+2})\big]$ with $f(x)=\arctan(\sqrt3x/(2+x))$, increasing on
$[0,1]$ with $f(1)=\pi/6$. Each bracket is positive, and $f(q^{3A+1})-f(q^{3A+2})\le f(q^{3A+1})-f(q^{3A+4})$ telescopes
to at most $f(q^{3n+1})<\pi/6$.
\end{proof}

\begin{lemma}[closed forms]\label{lem:c-closed}
Let $\chi'(r)=1$ for $r\equiv1$, $-1$ for $r\equiv2$, and $h_r(q)=(q^{Nr}+q^{(N+1)r})/(1-q^{3r})$. For $|q|<1$,
\[
\log E_n(\omega q)=S(q)+i\Delta(q),\qquad
S=-\sum_{3\nmid r}\frac{3}{2r}h_r+\sum_{3\mid r}\frac3r h_r,\qquad
\Delta=\sum_{3\nmid r}\frac{3\sqrt3}{2r}\chi'(r)\frac{q^{Nr}(1-q^r)}{1-q^{3r}},
\]
as an identity of analytic functions. Consequently $F(q):=\frac1{2i}\big(E_n(\omega q)-E_n(\bar\omega q)\big)=e^{S}\sin\Delta$
on $0<q<1$, and $F$ is the analytic continuation of this expression.
\end{lemma}
\begin{proof}
Expand $-3\log(1-\omega^kq^k)=3\sum_r\omega^{kr}q^{kr}/r$ and sum the geometric series over $k\equiv1,2\pmod3$, $k>3n$;
this gives $\sum_r\frac3r(\omega^rq^{Nr}+\omega^{2r}q^{(N+1)r})/(1-q^{3r})$, whose real and imaginary parts for real $q$
are the displayed expressions, and the displayed combination equals the series identically in $q$.
\end{proof}

Let $\chi(i)=1,-1,0$ for $i\equiv0,1,2\pmod 3$. With $B=\frac12$ and $C=2i-\frac12$ we have
$\cP_3(i)=\sqrt{a_3/(i-\frac14)}\,I_1\big(2\sqrt{a_3(i-\frac14)}\big)$.

\begin{theorem}[explicit expansion of $G_3^3$]\label{thm:expansion3}
For all $i\ge1$,
\[
s(i)=\sqrt3\,\chi(i)\,\cP_3(i)+\rho(i),\qquad |\rho(i)|\le\sum_{\substack{6\le k\le K(i)\\3\mid k}}\varphi(k)\cP_k(i)+\Econ,
\]
with $K(i)=\lfloor\sqrt{4\pi i}\rfloor+2$ and $\Econ=50.1501$.
\end{theorem}
The proof is in \S\ref{sec:expansion3}. Thus $\fa(i)=\sqrt3\chi(i)$ in the notation of~\eqref{eq:expansion-shape}.

\begin{lemma}[Poisson form]\label{lem:c-poisson}
Let $\cP^\sharp(q)=\sum_{i\ge1}\cP_3(i)q^i$. For $\Re w>0$, $|\Im w|\le\pi$,
\[
\cP^\sharp(e^{-w})=e^{-w/4}\big(e^{a_3/w}-1+\varepsilon(w)\big),\qquad
\varepsilon(w)=\sum_{k\ne0}(-i)^k\big(e^{a_3/(w+2\pi ik)}-1\big),
\]
and $|{-1}+\varepsilon(w)|\le 1.42$.
\end{lemma}
\begin{proof}
$x\mapsto\sqrt{a_3/x}\,I_1(2\sqrt{a_3x})$ has Laplace transform $e^{a_3/w}-1$; Poisson summation over the shifted lattice
$i-\frac14$ gives the identity. The bound is Certificate~\ref{cert:c-poisson}.
\end{proof}

\begin{certificate}\label{cert:c-poisson}
Ball-arithmetic enclosure of $|{-1}+\varepsilon(w)|$ on $1600$ boxes covering $0<\Re w\le10$, $|\Im w|\le\pi$, with the
terms $|k|>400$ bounded by $8a_3/(\pi\cdot400)+a_3^2e^{a_3/\pi}/(400\pi^2)$. Output: $|{-1}+\varepsilon|\le1.420$.
\cert{poisson\_check.py}{cubic/logs/04a\_poisson.log}
\end{certificate}

\begin{proposition}[reduction]\label{prop:c-reduce}
Let $m\equiv2\pmod3$, $W>0$ and $w=W+iy$. Then $c_n(m)=-2\Im J+R(m)$, where
\[
\Im J=\frac1{2\pi}\int_{-\pi}^{\pi}\Re\, e^{\Rh(w)}\,dy+\Im J_\varepsilon,\qquad
\Rh(w)=\frac{a_3}w+\Big(m-\frac14\Big)w+S(e^{-w})+\log\sin\Delta(e^{-w}),
\]
$|\Im J_\varepsilon|\le\frac1{2\pi}\int_{-\pi}^{\pi}1.42\,e^{(m-\frac14)W}|F(e^{-w})|\,dy$, and, for every $w_1>0$,
\[
|R(m)|\le\sum_{\substack{6\le k\le K(m)\\3\mid k}}3\varphi(k)\,e^{mw_1+a_k/w_1}E_n(e^{-w_1})+\Econ\,e^{mw_1}E_n(e^{-w_1}),
\qquad a_k=\pi^2/k^2 .
\]
\end{proposition}
\begin{proof}
$c_n(m)=\sum_je(j)s(m-j)$ with $e(j)=[q^j]E_n\ge0$. Inserting Theorem~\ref{thm:expansion3}, the $\chi$-part equals
$\sum_je(j)\sqrt3\chi(m-j)\cP_3(m-j)=-2\Im[q^m]\cP^\sharp(q)E_n(\omega q)$ for $m\equiv2$ (compare coefficients: $e(j)$ is
real and $\sqrt3\chi(m-j)=-2\Im\omega^{j}$ for $m\equiv2$). The coefficient is
$\frac1{2\pi}\int\cP^\sharp(e^{-w})E_n(\omega e^{-w})e^{mw}dy$; by Lemma~\ref{lem:c-poisson}, reality of
$[q^m]\cP^\sharp F$ and Lemma~\ref{lem:c-closed} this gives the displayed $\Im J$. For $R$, the remainder terms are bounded
termwise using $e(j)\ge0$, $\cP_k\ge0$ and $\sum_i\cP_k(i)e^{-w_1i}\le(1+2a_k)e^{a_k/w_1}\le3e^{a_k/w_1}$ (from
$I_1(y)\le\frac y2e^y$ and the Laplace transform).
\end{proof}

\begin{remark}\label{rem:c-identity}
Proposition~\ref{prop:c-reduce} is checked end to end by \texttt{identity\_check.py}: for $n=12$ ($m=200,302,404$) and
$n=20$ ($m=500,800,1100$) the exact $c_n(m)$ agrees with $-2(\Im J_0+\Im J_\varepsilon)+R(m)$, each piece computed
independently, to relative error $<10^{-8}$. The bound of Theorem~\ref{thm:expansion3} is checked exactly for all
$i\le3000$ and for $3388$ values $i\le108000$ (\texttt{expansion\_check.py}); the maximal ratio is $\sqrt3/2$. Neither
check is part of the proof.
\end{remark}

In particular, $c_n(m)<0$ follows from
\begin{equation}\label{eq:c-target}
\frac1{2\pi}\int_{-\pi}^{\pi}\Re\,e^{\Rh(W+iy)}\,dy\ >\ |\Im J_\varepsilon|+\tfrac12|R(m)| .
\end{equation}

\subsection{The finite range and the band}\label{sec:c-finite}

\begin{certificate}[finite range]\label{cert:c-finite}
$P_n^3$ is computed exactly in $\Z[q]$ (FLINT) for $n\le243$, and it is checked that $c_n(m)=0$ for $m<3n$ and $c_n(m)<0$
for $3n\le m\le\frac92n^2$, $m\equiv2$. Output: no violations.
\cert{exact\_small\_n.py}{cubic/logs/01\_exact\_small\_n\_\{a,b,c\}.log}
\end{certificate}

\begin{proposition}[band]\label{prop:c-band}
For $3n+2\le m<6n+2$ with $m=3M+2$, $c_n(m)=3\sum_{i=0}^{M-n}g(i)$ with $g(i)=s(3i)+s(3i+1)$, and $g(i)<0$ for all
$i\ge0$. Hence $c_n(m)<0$ in the band.
\end{proposition}
\begin{proof}
Since $6n+2=2N$, $E_n\equiv1+3\sum_{N\le j<2N,\,3\nmid j}q^j\pmod{q^{2N}}$, and $s(m)=0$. Grouping $j=3A+1$ and $j=3A+2$
gives the identity (checked exactly for $n\le60$ by \texttt{band\_lemma.py}). Negativity of $g$ is exact for $i\le36000$
(Certificate~\ref{cert:c-band}). For $i\ge36000$, Theorem~\ref{thm:expansion3} gives
$g(i)\le-\sqrt3\big(\cP_3(3i+1)-\cP_3(3i)\big)+\rho_{\max}(3i)+\rho_{\max}(3i+1)$. Since $\cP_3(x)=2a_3I_1(y)/y$ with
$y=2\sqrt{a_3(x-\frac14)}$, we have $\cP_3'(x)=4a_3^2I_2(y)/y^2$, and $\cP_3$ is increasing and convex, so the difference
is at least $\cP_3'(3i)$. Using $I_2(y)\ge\kappa_2e^y/\sqrt{2\pi y}$ for $y\ge y_0$, $I_1(y)\le e^y/\sqrt{2\pi y}$,
$\cP_k\le\cP_6$ for $k\ge6$ and $\sum\varphi(k)\le K^2/2$, the ratio of the $\rho$-bound to the main difference is of the
form ${\rm poly}(y)e^{-y/2}$, decreasing for $y\ge9$, and at $i=36000$ ($y_0=688.3$) it is at most $5\cdot10^{-139}$ in
ball arithmetic.
\cert{band\_tail\_rigorous.py}{cubic/logs/02b\_band\_tail.log}
\end{proof}

\begin{certificate}\label{cert:c-band}
Exact $s(i)$ from Jacobi's identity and the recurrence $jb(j)=3\sum_t\sigma(t)b(j-t)$ for $b=[q^j](q;q)^{-3}_\infty$.
Output: the identity of Proposition~\ref{prop:c-band} holds for $n\le60$; $g(i)<0$ for all $i\le36000$.
\cert{band\_lemma.py 36000}{cubic/logs/02\_band\_lemma.log}
\end{certificate}

\subsection{The Laplace regime}\label{sec:c-laplace}

\subsubsection{Scaled form}
Let $\zeta=Nw$, $\epsilon=1/N$, $g(x)=x/(1-e^{-x})$ and $h(x)=(1-e^{-x})/(1-e^{-3x})=g(3x)/(3g(x))$. Both $g$ and $h$ are
analytic for $|x|<2\pi/3$.

\begin{lemma}\label{lem:c-scaled}
$S(e^{-w})=N\St(\zeta,\epsilon)$ and $\Rh(w)=N\,\cG(\zeta)+(m-N-\frac14)\zeta/N$, where
\[
\St=\sum_r s_r\,e^{-r\zeta}\big(1+e^{-r\zeta\epsilon}\big)\frac{g(3r\zeta\epsilon)}{3r\zeta},\qquad
\Qt=\sum_{3\nmid r}\frac{3\sqrt3}{2r}\chi'(r)e^{-(r-1)\zeta}h(r\zeta\epsilon),
\]
$s_r=-\frac3{2r}$ ($3\nmid r$), $\frac3r$ ($3\mid r$), and
$\cG=\frac{a_3}\zeta+\St+\epsilon\big(\log\Qt+\log\tfrac{\sin\Delta}\Delta\big)$ with $\Delta=\Qt e^{-\zeta}$. The saddle
condition $\Rh'(W)=0$ reads $-\cG'(v)=\tilde\mu$ with $v=NW$.
\end{lemma}
\begin{proof}
Substitute $q=e^{-\zeta\epsilon}$ in Lemma~\ref{lem:c-closed}; $\log\sin\Delta=\log\Qt-\zeta+\log\frac{\sin\Delta}{\Delta}$,
and the $-\zeta$ combines with $(m-\frac14)w$ to $(m-N-\frac14)\zeta/N$.
\end{proof}

Thus a pair $(N,m)$ in this regime is represented by its real saddle $v$ and $\epsilon=1/N$. Since $m\ge2N$,
$\tilde\mu\ge\epsilon(1-\frac1{4N})\ge0.999\epsilon$; balls with $-\cG'(v)<0.999\epsilon$ on the whole ball contain no pairs.

\subsubsection{The error budget on a ball}
Fix a ball $V_b\ni v$ and a ball $E\ni\epsilon$ with $E\subset[0,1/733]$, and let $N_b=\max(733,1/\sup E)$. Write
$\cG_j=\cG^{(j)}(v)$, $\sigma_y=N^{-3/2}\cG_2^{-1/2}$ and $\tau=y/\sigma_y$. Dividing~\eqref{eq:c-target} by
$e^{\Rh(W)}\sigma_y/(2\pi)$, it suffices to show
\begin{equation}\label{eq:c-B}
B_{\rm tot}:=\frac{Z_1+Z_2+Z_3+Z_\varepsilon+Z_R}{L_{\rm w}}<1,
\end{equation}
where $L_{\rm w}$ bounds the window integral from below and the $Z$'s bound the remaining contributions from above.

\smallskip\noindent\emph{Window $|\tau|\le c$.}
$\Rh(W+i\sigma_y\tau)-\Rh(W)=-\frac{\tau^2}2-i\cG_{3,n}\frac{\tau^3}6+\phi_4\frac{\tau^4}{24}+E_5$ with
$\cG_{3,n}=\sqrt\epsilon\,|\cG_3|\cG_2^{-3/2}$, $\phi_4=\epsilon \cG_4\cG_2^{-2}$, $|E_5|\le \cG_{5,n}|\tau|^5/120$,
$\cG_{5,n}=\epsilon^{3/2}M_5\cG_2^{-5/2}$, where $M_5$ bounds $|\cG^{(5)}|$ on the window: $120a_3/v^6$ for the $a_3/\zeta$
part plus a Cauchy estimate (Lemma~\ref{lem:cauchy}) for the rest on the circle $|\zeta-v|=0.6v$.
$L_{\rm w}=\int_{-c}^{c}\ell(\tau)d\tau$ with $\ell$ the resulting pointwise lower bound of the real part
($\cos\theta\ge1-\theta^2/2$).

\smallskip\noindent\emph{Zone 1, $c\le|\tau|$, $|t|\le t_s$} ($t=Ny$). $\Re[\cG(v+it)-\cG(v)]\le-\frac{t^2}2(\kappa_0-\epsilon\kappa_b)$
from the even Taylor terms of $\cG_0=a_3/\zeta+\St$ and of $b=\log\Qt+\log\frac{\sin\Delta}\Delta$ with Cauchy remainders at
order $6$; thus $Z_1=2\int_c^\infty e^{-\kappa\tau^2/2}d\tau$, $\kappa=(\kappa_0-\epsilon\kappa_b)/\cG_2$, independent of $N$.

\smallskip\noindent\emph{Zone 2, $t_s\le t\le t_C$.} On cells, $\Re\Rh-\Rh(W)=NA(t)+B(t)$ with
$A=\Re[\cG_0(v+it)-\cG_0(v)]$ and $B=\Re[b(v+it)-b(v)]$, enclosed in ball arithmetic (the $a_3/\zeta$ part in closed form,
the factor $e^{-rv}$ separated to control dependency). The certificate asserts $A<-1/(2N_b)$ on every cell, so
$N^{1/2}e^{NA}\le N_b^{1/2}e^{N_bA}$ for $N\ge N_b$, and
$Z_2=(N_b\cG_2)^{1/2}\sum_{\rm cells}2(e^{N_bA+B}+3e^{N_b(A-\Re(a_3/\zeta))+B})\,|{\rm cell}|$ (the second term is the
$\varepsilon$-part).

\smallskip\noindent\emph{Zone 3, $t\ge t_C$.} Either (crude) $|F|\le E_n(e^{-W})$ and
$|e^{a_3/w}-1+\varepsilon|\le e^{a_3W/(W^2+y^2)}+3$, giving an exponent $N\beta+O(1)$ with
$\beta=\Psi(v)/v+a_3v/(v^2+t_C^2)-a_3/v-\St(v)$, $\Psi(v)=2e^{-v}/(1-e^{-v})$, asserted $\beta<-3/(2N_b)$; or (split,
when $\inf E>0$) $|\Qt|\le\sum_{3\nmid r}\frac{3\sqrt3}{2r}e^{-(r-1)v}2(1/(3rv\inf E)+1)$ and $|\St|$ bounded termwise,
with slope $<-5/(2N_b)$.

\smallskip\noindent\emph{$R(m)$.} Proposition~\ref{prop:c-reduce} with $w_1=t'W$, $t'$ from a fixed list,
$\log E_n(e^{-w_1})\le N\Psi(Nw_1)/(Nw_1)+6|\log(1-e^{-Nw_1})|$, $\tilde\mu=-\cG'(v)$ and $\varphi$-sum $\le3(K^2/6+K)$;
each exponent is asserted to have slope $<-7/(2N_b)$.

\begin{lemma}\label{lem:c-monotone}
If all assertions above hold on $(V_b,E)$, then $B_{\rm tot}$ computed with $N=N_b$ bounds~\eqref{eq:c-B} for every $N$
with $1/N\in E$ and every $m$ whose saddle lies in $V_b$.
\end{lemma}
\begin{proof}
$L_{\rm w}$ and $Z_1$ do not depend on $N$ (their $N$-dependence enters only through $\epsilon\in E$). Every other term is
of the form $N^pe^{N\lambda}$ with $\lambda$ bounded by the asserted slope $<-p/N_b$, hence non-increasing in $N\ge N_b$.
\end{proof}

\begin{certificate}[uniform Laplace grid]\label{cert:c-laplace}
$v\in[2.95,3000]$ is covered by $36$ chunks of balls and, for each $v$-ball, $\epsilon\in[0,1/733]$ by $\epsilon$-balls;
balls whose pairs are vacuous are discarded only after the rigorous bound
$\tilde\mu\le a_3/v^2+\max_{|\zeta-v|=v/2}|h(\zeta)-h(v)|/(v/2)$ (with $h=\St+\epsilon b$) proves
$\tilde\mu<0.999\epsilon$, or when $\tilde\mu>0.0985$ (centre regime). Failing balls are bisected. Output: every chunk
certified, $0$ unresolved balls, $\max B_{\rm tot}=0.117$.
\cert{run\_uniform.py}{cubic/logs/uni\_*.log (36 logs)}
\end{certificate}

\subsubsection{The closing lemma}\label{sec:c-closing}

\begin{proposition}\label{prop:c-closing}
For every pair with saddle $v\ge3000$, \eqref{eq:c-B} holds with $B_{\rm tot}\le0.0082$.
\end{proposition}
\begin{proof}[Proof and certificate]
Put $s=\epsilon v^2/a_3$; the constraint $\tilde\mu\ge0.999\epsilon$ and $\tilde\mu=(a_3/v^2)(1+O(e^{-v}))$ give
$s\le1.0011$. On $|\zeta-v|\le v/2$: $|\St|\le\sum_r\frac3re^{-rv/2}2(\frac2{3rv}+\epsilon)\le5\cdot10^{-655}$; $|h(x)|\le4$
for every $r$ (by the Bernoulli enclosure if $|x|<0.9$, and $|h(x)|\le2(|x|/(3\Re x)+1)$ with
$|x|/\Re x=|\zeta|/\Re\zeta\le3$ otherwise), so the $r\ge2$ part of $\Qt$ and $\Delta$ are $\le4\cdot10^{-651}$; and
$\log h$ on the disc $|x|\le\frac32a_3s/v$ is enclosed, giving $|b(\zeta)-b(v)|\le4.4\cdot10^{-3}$. By
Lemma~\ref{lem:cauchy}, the relative perturbations of $\cG_2,\dots,\cG_5$ from their pure $a_3/\zeta$ values are
$\le4.7\cdot10^{-5}$, so $\cG_{3,n}\le0.0388$, $\phi_4\le2.0\cdot10^{-3}$, $\cG_{5,n}\le1.03\cdot10^{-3}$,
$\kappa\ge0.8(1-10^{-5})$ on $|t|\le v/2$, and $L_{\rm w}\ge2.494$, $Z_1\le0.0205$. For $|t|\ge v/2$ the exponent is at most
$-\frac{v}{5s}(1-10^{-5})+\log(\frac{4v}{3a_3s}+6)+20ve^{-v}/(a_3s)$ with prefactor $\sqrt2v^{3/2}/(a_3s^{3/2})$; its
$s$-derivative is positive for $s\le1.0011$ because $v/5>2.5\cdot1.0011+20ve^{-v}/a_3$, and it decreases in $v$, so the
worst case is $v=3000$, $s=1.0011$, where $Z_2\le5\cdot10^{-251}$. The $R$-term is $\le10^{-1278}$. All numbers are ball
enclosures.
\cert{ray\_asymptotic2.py}{cubic/logs/07\_closing\_lemma.log}
\end{proof}

\subsection{The centre regime}\label{sec:c-centre}

Here $c_n(m)=\frac1{2\pi}\int_{-\pi}^{\pi}P_n(re^{i\theta})^3r^{-m}e^{-im\theta}d\theta$ is treated directly, without
Proposition~\ref{prop:c-reduce}. Put $\cL(u)=3\sum_{k\le3n,3\nmid k}\log(1-e^{ku})$, $u_0=2\pi i/3$, and let $u^*=u_0+x^*$
be the saddle $\cL'(u^*)=m$ near $u_0$; write $\xi=-N\Re x^*$ and $\eta=N\Im x^*$. The conjugate saddle near $-u_0$
contributes the complex conjugate.

\begin{lemma}[sign margin]\label{lem:c-phase}
The phase of the main term is $\alpha(r)-\frac{2\pi}3+b+E_1\pmod{2\pi}$, where
$\alpha(r)=-\frac\pi6+\arctan\big(\sqrt3r^{3n}/(2+r^{3n})\big)$, $b=-\frac12\arg \cL''(u^*)$, and
$|E_1|\le\sup|\cL''|\,\delta^2/2$ with $\delta=\Im x^*$; moreover
$|\delta|\le3N\TV_s/(N^3\rho_{\rm lo}-3N^2V_{\rm K})$ by the saddle equation and Abel summation.
\end{lemma}
\begin{proof}
The main term is $M=e^{\cL(u^*)-mu^*}\,\omega\,e^{ib}\sqrt{2\pi/|\cL''(u^*)|}/(2\pi)$ (steepest descent through $u^*$ in the
direction $ie^{ib}$). Write $u_0'=\log r+2\pi i/3$, so $u^*=u_0'+i\delta$. Then
$\Im[\cL(u^*)-mu^*]=\Im \cL(u_0')+\int_0^\delta(\Re \cL'(u_0'+is)-m)\,ds-m\cdot\frac{2\pi}3$, and $\Re \cL'(u^*)=m$ gives
$|\Re \cL'(u_0'+is)-m|\le\sup|\cL''|\,|\delta-s|$, whence the bound on $E_1$. Next $\Im \cL(u_0')=\arg P_n(r\omega)^3=\alpha(r)$
by the pairing $(1-r^{3j+1}\omega)(1-r^{3j+2}\omega^2)$ of Lemma~\ref{lem:c-tailphase} type; $m\equiv2$ gives
$m\cdot\frac{2\pi}3\equiv\frac{4\pi}3$, and $\arg\omega=\frac{2\pi}3$. For $\delta$:
$\Im \cL'(u_0')=3\sum_i[(3i+2)s_\omega(x_{3i+2})-(3i+1)s_\omega(x_{3i+1})]$ with $s_\omega(x)=\Im(\omega x/(1-\omega x))$ and
$x_k=r^k$, which is at most $3N\,\TV_\tau(\tau s_\omega(e^{-\xi\tau}))$ in modulus; and $\Im \cL'(u_0'+i\delta)=0$ with
$\Re \cL''\ge N^3\rho_{\rm lo}-3N^2V_{\rm K}$ on the segment.
\end{proof}
(Here $\rho_{\rm lo}$ is a lower bound for the scaled continuum curvature and $V_{\rm K}$ a Koksma total-variation constant,
Lemma~\ref{lem:koksma}.)

The certificate splits the circle into a window $|\theta-\theta^*|\le c\sigma$, a middle stretch up to
$|\theta-2\pi/3|\le0.05$, and the far region. Each contribution is bounded uniformly in $N\ge733$ as a function of $\xi$:
\begin{itemize}
\item \emph{Window}: the cubic Taylor constant $G_3=H_3\sigma^3$ with
$H_3(\text{window})\le N^4[h_3+3(\TV+S)/N]$ and $\sigma^3\le(N^3a_2-3N^2V_{\rm K})^{-3/2}$ (continuum integrals plus Koksma
remainders, Lemma~\ref{lem:koksma}); the window shrinks with $N$.
\item \emph{Middle}: three zones in $t=N(\theta-2\pi/3)$: Taylor with $\Re \cL''\ge N^3\rho_{\rm lo}(t)-3N^2V_{\rm K}$
($|t|\le0.8$); the linearised continuum profile $-N\int_0^1 g(e^{-\xi\tau})(1-\cos\tau t)d\tau$ plus a Koksma constant
($0.8\le t\le10$, slope asserted); a Dirichlet-kernel bound ($t\ge10$, slope asserted).
\item \emph{Far}: for $\theta\ge0.05$ a Dirichlet-kernel bound; near $\theta=0$ a continuum profile
$-N\int g(1+2\cos\tau t)$ with positivity asserted on $t\in[0,40]$ and a kernel bound beyond; slopes asserted.
\item \emph{Sign margin}: Lemma~\ref{lem:c-phase} with $r^{3n}=e^{-\xi(1-1/N)}$ enclosed over all $N\ge733$.
\end{itemize}

\begin{certificate}[centre tables]\label{cert:c-centre}
On $74$ $\xi$-balls (with finer sub-tables of $115$ balls for the middle and window constants) covering $[-0.3,3.4]$, the
programs produce uniform bounds for the middle stretch, the window constant, the sign margin (with $\cos<0$ asserted) and
the far region. \texttt{coverage.py} (check E) combines them into
${\rm err}=(2(\text{inner}+\text{middle})+\text{far})/(2\,|\cos|)$ on every ball. Output: $\max{\rm err}=0.49<1$. The
saddle offset satisfies $|\eta|\le0.0155$, within the shift margins of the tables (check G).
\cert{centre\_mid\_uniform\{,2,3\}.py, centre\_g3\_uniform\{,2,3\}.py, centre\_clow\_uniform.py, coverage.py}{cubic/tables/*\_table.txt, cubic/logs/08\_coverage.log}
\end{certificate}

\subsection{Coverage and proof of Theorem~\ref{thm:cubic-main}}\label{sec:c-coverage}

\begin{proposition}\label{prop:c-coverage}
Every $(n,m)$ with $n\ge244$, $m\equiv2\pmod3$, $6n+2\le m\le\frac92n^2$ lies in the Laplace regime with saddle in
$[2.95,3000]$ or $[3000,\infty)$, or in the centre regime with $\xi\in[-0.3,3.4]$.
\end{proposition}
\begin{proof}
(A) The chunks of Certificate~\ref{cert:c-laplace} tile $[2.95,3000]$ with no unresolved balls. (B) $-\cG'(2.95)\ge0.1015>0.0985$
for all $\epsilon\in[0,1/733]$ and $\cG''>0$, so every pair with $\tilde\mu<0.0985$ has $v>2.95$. (C) $m\ge2N$ gives
$\tilde\mu\ge0.999\epsilon$. (D) $\tilde\mu\ge0.0985$ implies $\bar\mu>\tilde\mu\ge0.098$. (F)
$\bar\mu_N(\xi)=N^{-2}\Re \cL'(u^*)$ is decreasing in $\xi$ and, by Lemma~\ref{lem:koksma}, $\bar\mu_N(3.4)\le0.0886<0.098$
and $\bar\mu_N(-0.3)\ge0.531>\frac12$ for all $N\ge733$, while $m\le\frac92n^2$ gives $\bar\mu<\frac12$. Each item is checked
by \texttt{coverage.py}.
\cert{coverage.py}{cubic/logs/08\_coverage.log}
\end{proof}

\begin{proof}[Proof of Theorem~\ref{thm:cubic-main}]
Certificate~\ref{cert:c-finite} for $n\le243$; for $n\ge244$ Proposition~\ref{prop:c-band} for $m<6n+2$, and otherwise
Proposition~\ref{prop:c-coverage} with Lemma~\ref{lem:c-monotone}, Certificate~\ref{cert:c-laplace},
Proposition~\ref{prop:c-closing} and Certificate~\ref{cert:c-centre}.
\end{proof}

\medskip\noindent\emph{Load-bearing hand arguments.} The proof of Theorem~\ref{thm:cubic-main} depends on the following
written arguments in addition to the certificates: Lemmas~\ref{lem:c-tailphase}--\ref{lem:c-poisson},
Proposition~\ref{prop:c-reduce}, Theorem~\ref{thm:expansion3} (\S\ref{sec:expansion3}), Lemma~\ref{lem:c-scaled}, the
derivation of the budget~\eqref{eq:c-B}, Lemma~\ref{lem:c-monotone}, the analytic parts of
Proposition~\ref{prop:c-closing}, Lemma~\ref{lem:c-phase}, and the elementary lemmas of Appendix~\ref{app:elementary}.

\subsection{Proof of Theorem~\ref{thm:expansion3}}\label{sec:expansion3}

We specialise~\eqref{eq:Gp-transform} to $p=\delta=3$; then $G_3^3=P_\infty^3$. Let $\tau=h/k+iz/k^2$, $d=\gcd(3,k)$.
The eta transformation gives
\[
G_3(e^{2\pi i\tau})^3=\Big(\frac3d\Big)^{3/2}e^{3\pi i(-s(h,k)+s(3h/d,k/d))}
\exp\Big(\frac{\pi(d^2-3)}{12z}-\frac{\pi z}{2k^2}\Big)\widehat G^3,\qquad
\widehat G=\frac{f(e^{2\pi i dh'_d/k-2\pi d^2/(3z)})}{f(e^{2\pi ih'/k-2\pi/z})},
\]
with $hh'\equiv-1\pmod k$ and principal powers.

\emph{Growing and bounded cusps.} If $3\mid k$ then $d=3$, the prefactor has modulus $1$ and the exponential is
$e^{\pi/(2z)}$: every such cusp grows with $B=\frac12$, and the correction term of $\widehat G^3$ grows only if
$\delta>24/(p-1)=12$, which is not the case. If $3\nmid k$, then on $\Re(1/z)=w\ge1$, $|G_3^3|\le K_3^3$ with
$K_3=\sqrt3\,e^{-\pi/18}f(e^{-2\pi/3})f(e^{-2\pi})=1.691256$, the maximum being attained at $w=1$.

\emph{Rademacher dissection.} Take Farey order $N'\ge\sqrt{4\pi i}+1$. On the Ford-circle arcs mapped to
$K=\{|z-\frac12|=\frac12\}$, the endpoints satisfy $|z'|,|z''|\le\sqrt2k/(N'+1)$ and $\Re z\le2k^2/(N'+1)^2$, so
$|e^{\pi Cz/k^2}|\le e$ with $C=2i-\frac12$.

\emph{Main terms.} For $3\mid k$, the term $1$ of $\widehat G^3$ gives $\exp(\pi/(2z)+\pi Cz/k^2)$; completing the arc to
$K$ and using $\frac i{k^2}\int_K\exp(\pi B/z+\pi Cz/k^2)dz=\cP_k(i)$ (Lemma~\ref{lem:bessel-contour}) yields the main term
with amplitude of modulus $1$, for each of the $\varphi(k)$ values of $h$. The completion costs at most
$\frac1{k^2}\cdot2\cdot\frac\pi2\cdot\frac{\sqrt2k}{N'+1}e^{1+\pi/2}$ per cusp, and
$\sum_{3\mid k\le N'}\varphi(k)/k\le\frac29(N'+1)$ because $\varphi(3j)/(3j)\le\frac23$; total
$E_{\rm arc}=\sqrt2\pi\frac29e^{1+\pi/2}=12.9103$.

\emph{Remainder at growing cusps.} Moving to the chord (length $\le2\sqrt2k/(N'+1)$),
$|e^{\pi/(2z)}(\widehat G^3-1)|\le e^{\pi w/2}(f(e^{-2\pi w})^3f(e^{-6\pi w})^3-1)$, whose terms $e^{\pi w/2}e^{-2\pi jw}$
($j\ge1$) are non-increasing, so the maximum is at $w=1$; total
$E_{\rm rem}=2\sqrt2\frac29e^{1+\pi/2}(f(t)^3f(t^3)^3-1)=0.0463$, $t=e^{-2\pi}$.

\emph{Bounded cusps.} On the chords, the integrand is at most $eK_3^3$; summing $\frac1{k^2}\frac{2\sqrt2k}{N'+1}$ over
$\le\varphi(k)\le k$ values of $h$ and $k\le N'$ gives $E_{\rm ng}=2\sqrt2eK_3^3=37.1935$.

\emph{The $k=3$ amplitude.} For $k=3$ we have $d=3$, $(3/d)^{3/2}=1$ and $s(3h/d,k/d)=s(h,1)=0$, so the phase of the cusp
$h/3$ is $e^{-3\pi is(h,3)}$. The Dedekind sums are $s(1,3)=\frac1{18}$ and $s(2,3)=-\frac1{18}$, and the leading term of
$\widehat G^3$ is $1$. Hence the two cusps contribute $\cP_3(i)$ times
\[
e\big(-\tfrac i3\big)e^{-\pi i/6}+e\big(-\tfrac{2i}3\big)e^{\pi i/6}=2\cos\Big(\frac{2\pi i}3+\frac\pi6\Big)=\sqrt3,\ -\sqrt3,\ 0
\quad(i\equiv0,1,2\bmod 3),
\]
that is, $\sqrt3\chi(i)\cP_3(i)$, where $e(x)=e^{2\pi ix}$. (This is also confirmed by the exact agreement recorded in
Remark~\ref{rem:c-identity}.) For $k\ge6$ only the modulus $1$ of each amplitude is used. Summing,
$\Econ=E_{\rm arc}+E_{\rm rem}+E_{\rm ng}=50.1501$.
\cert{sz\_expansion\_audit.py (recomputes the constants)}{cubic/logs/03\_sz\_audit.log}
\qed

%=====================================================================
\section{The Third Borwein conjecture}\label{sec:third}
%=====================================================================

In this section $p=5$, $\delta=1$, $P_n=P_{5,n}=(q;q)_{5n}/(q^5;q^5)_n$, $c_n(m)=c(m)=[q^m]P_n$, $\deg P_n=10n^2$,
$N=5n+1$, $a_5=2\pi^2/75$, $\omega_5=e^{2\pi i/5}$, $s(i)=[q^i]G_5$, $E_n=\prod_{k>5n,\,5\nmid k}(1-q^k)^{-1}$, $\lambda=m/N$
and $\mu=m/(10n^2)$. The amplitude of the $k=5$ main term is
\[
\fa(i)=2\cos(2\pi i/5+\pi/5)+2\cos(4\pi i/5)=\tfrac{5+\sqrt5}2,\ -\sqrt5,\ -\tfrac{5-\sqrt5}2,\ 0,\ 0\qquad(i\equiv0,1,2,3,4).
\]
It vanishes on the thin classes $3,4$ of Lemma~\ref{lem:vanish}(b). In the classes $0,1,2$, $s(m)$ has the non-vanishing
main term $\fa(m)\cP_5(m)$ (Theorem~\ref{thm:G5}).

\begin{theorem}\label{thm:third-main}
For every $n\ge1$ and every $m$, $c_n(m)\ge0$ if $5\mid m$ and $c_n(m)\le0$ otherwise. That is, the Third Borwein
conjecture holds.
\end{theorem}

Labels. Each step below is of one of the following kinds: a complete written argument; an inequality verified in ball
arithmetic over its whole parameter range (``certified''; a \emph{Certificate} line names the program and log); a
verification by exact integer arithmetic (``exact''); or a standard fact (Appendix~\ref{app:facts}).

\subsection{Statement, reduction and regions}\label{sec:t-reduction}

By Lemma~\ref{lem:palin} it suffices to prove the sign pattern for $0\le m\le5n^2$, that is, for $0\le\mu\le\frac12$. The
three analytic regions are
\begin{itemize}
\item \textbf{band:} $\lambda<4$ (that is, $m<4N$);
\item \textbf{Laplace region:} $\lambda\ge4$ and $\mu<0.105$;
\item \textbf{centre:} $\mu\ge0.105$.
\end{itemize}
Their union is $\{0\le m\le5n^2\}$ trivially. Table~\ref{tab:regions5} summarises the covering;
Theorem~\ref{thm:t-coverage} checks that each region's certificate covers the whole of its region for $n\ge980$.

\begin{table}[h]
\centering
\small
\begin{tabular}{llll}
\toprule
region & classes & method & where\\
\midrule
$n\le979$ & all & exact integer computation & \S\ref{sec:t-exact}\\
$\lambda<4$ & $3,4$ & one-part block, exact to $n\le2000$, majorant $\Psi$ beyond & \S\ref{sec:t-band34}\\
$\lambda<4$ & $0,1,2$ & exact to $n\le2000$, per-$m$ inequality beyond & \S\ref{sec:t-band012}\\
$\lambda\ge4$, $\mu<0.105$ & $0,1,2$ & saddle point, $V\in[2,12]$ tiles and tail cell & \S\ref{sec:lap012}\\
$\lambda\ge4$, $\mu<0.105$ & $3,4$ & re-centred (thin) saddle, tiles and tail cell & \S\ref{sec:lap34}\\
$\mu\ge0.105$ & all & near-peak Laplace, off-peak lemma, combination & \S\ref{sec:t-centre}\\
\bottomrule
\end{tabular}
\caption{The covering for $p=5$ (for $n\ge980$ in the analytic regions).}
\label{tab:regions5}
\end{table}

By Lemma~\ref{lem:factor}, $e(j)$ counts the partitions of $j$ into parts $>5n$ that are not divisible by $5$; for $j<4N$
at most three parts fit.

\begin{remark}[cusp by cusp]
In the expansion of \S\ref{sec:G5}, the $k=10$ and $k=15$ amplitudes also vanish on $i\equiv3,4\pmod 5$, to $390$ or more
digits (\texttt{cusp10.py}, \texttt{cusp15b.py}). This is not used in the proof.
\end{remark}

%---------------------------------------------------------------------
\subsection{The expansion of \texorpdfstring{$G_5$}{G5}}\label{sec:G5}
%---------------------------------------------------------------------

\emph{Notation for this subsection only.} The coefficient index is written $\nu$, and $i=\sqrt{-1}$. $N$ is the Farey
order, not $5n+1$. $e(x)=e^{2\pi ix}$; $f(x)=F(x)=1/(x;x)_\infty$; $s(h,k)=\sum_{r=1}^{k-1}\frac rk\big(\big(\frac{hr}k\big)\big)$
is the Dedekind sum, and $\omega(h,k)=e^{\pi is(h,k)}$; $t_0=e^{-2\pi}$.

\subsubsection{Statement}
Put $B=1/3$, $C=C(\nu)=2\nu-1/3$, $K(\nu)=\lfloor\sqrt{4\pi\nu}\rfloor+2$,
$\cP_k(\nu)=\frac{2\pi}k\sqrt{B/C}\,I_1\big(\frac{2\pi}k\sqrt{BC}\big)$ and $\fa(\nu)=\sum_{h=1}^4w_h\omega_5^{-\nu h}$,
where $w_h=e^{-\pi is(h,5)}$.

\begin{theorem}[explicit expansion of $G_5$]\label{thm:G5}
For every $\nu\ge1$,
\[
\big|s(\nu)-\fa(\nu)\cP_5(\nu)\big|\le\sum_{5\mid k,\ 10\le k\le K(\nu)}\varphi(k)\cP_k(\nu)+E_{\rm arc}+E_{\rm rem}+E_{\rm ng},
\]
where
\begin{itemize}
\item $E_{\rm arc}=\sqrt2\pi\cdot\frac4{25}e^{1+\pi/3}\le 5.506447$;
\item $E_{\rm rem}=2\sqrt2\cdot\frac4{25}e^{1+\pi/3}\big(f(t_0)f(t_0^5)-1\big)\le0.006571$, with $t_0=e^{-2\pi}$;
\item $E_{\rm ng}=2\sqrt2\,e\,K_5\le 21.950524$, with $K_5=\sqrt5e^{-\pi/15}f(e^{-2\pi/5})f(e^{-2\pi})\le2.8549955$.
\end{itemize}
Their sum is $\Econ\le 27.463542\le27.4636$.
\end{theorem}

Every $\cP_k(\nu)$ is positive: $C\ge5/3$, and $I_1>0$ on $(0,\infty)$. By \S\ref{sec:G5-constants},
$s(1,5)=1/5$, $s(2,5)=s(3,5)=0$ and $s(4,5)=-1/5$. So $w_1=e^{-\pi i/5}$, $w_2=w_3=1$, $w_4=e^{\pi i/5}$, and
$\fa(\nu)=2\cos(2\pi\nu/5+\pi/5)+2\cos(4\pi\nu/5)$, as above.

\subsubsection{Farey order}
Take $N=K(\nu)$. Two properties are used.
\begin{itemize}
\item \textbf{(N1)} $N\ge5$ for all $\nu\ge1$, since $\sqrt{4\pi}>3$. So the denominator $k=5$ always occurs. (The choice
$N=K(\nu)-1$ fails at $\nu=1$: it gives $N=4$.)
\item \textbf{(N2)} $N+1=\lfloor\sqrt{4\pi\nu}\rfloor+3>\sqrt{4\pi\nu}$. Hence $2\pi C/(N+1)^2=(4\pi\nu-2\pi/3)/(N+1)^2<1$.
\end{itemize}

\subsubsection{The Rademacher path}

\begin{lemma}[Farey dissection; Fact~\ref{fact:farey}]\label{lem:farey}
Take the Farey fractions of order $N$ cyclically on $[0,1)$, and let $h_1/k_1<h/k<h_2/k_2$ be neighbours. Let $C(h/k)$ be
the Ford circle with centre $h/k+i/(2k^2)$ and radius $1/(2k^2)$, and let $\gamma_{h,k}$ be its upper arc between its
tangency points with $C(h_1/k_1)$ and $C(h_2/k_2)$. Then these arcs join into a path from $\tau_0$ to $\tau_0+1$ in $\HH$,
and $s(\nu)=\int_{\tau_0}^{\tau_0+1}G_5(e(\tau))e(-\nu\tau)\,d\tau$, taken over a horizontal segment, may be deformed onto
this path by Cauchy's theorem. The integrand is holomorphic on $\HH$ and $1$-periodic.
\end{lemma}

The substitution $\tau=h/k+iz/k^2$ maps $C(h/k)$ onto $K:\ |z-\tfrac12|=\tfrac12$. Also $d\tau=(i/k^2)dz$ and
$e(-\nu\tau)=e(-\nu h/k)e^{2\pi\nu z/k^2}$. This gives
\begin{equation}\label{eq:G5-dissection}
s(\nu)=\sum_{k\le N}\ \sum_{\substack{0\le h<k\\(h,k)=1}}\frac i{k^2}e(-\nu h/k)\int_{z'}^{z''}G_5\big(e(h/k+iz/k^2)\big)e^{2\pi\nu z/k^2}dz,
\end{equation}
with $z'=(k^2+ikk_1)/(k^2+k_1^2)$ and $z''=(k^2-ikk_2)/(k^2+k_2^2)$. The integral runs clockwise along $K$ from $z'$
($\Im>0$) through $z=1$ to $z''$ ($\Im<0$). The tangency point of $C(h/k)$ with $C(h_1/k_1)$ is
$h/k-k_1/(k(k^2+k_1^2))+i/(k^2+k_1^2)$, and it maps to $z'$; the same holds for $z''$.

\emph{The arc of $0/1$.} In the cyclic dissection of $[0,1)$, the circle $C(0/1)$ contributes two pieces: the arc from
$\tau_0=i$ to its tangency with $C(1/N)$, and the arc of $C(1/1)=C(0/1)+1$ from its tangency with $C((N-1)/N)$ to
$\tau_0+1=i+1$. The integrand $G_5(e(\tau))e(-\nu\tau)$ is $1$-periodic. Translate the second piece by $-1$: it becomes the
arc of $C(0/1)$ from the tangency with $C(-1/N)$ to $i$. The two pieces then join into a single arc of $C(0/1)$, from the
tangency with $C(-1/N)$ through $i$ to the tangency with $C(1/N)$. In the $z$-variable ($h=0$, $k=1$, $\tau=iz$), this is
the arc of $K$ from $z'$ (with $k_1=N$) through $z=1$ (that is, $\tau=i$) to $z''$ (with $k_2=N$). So the term $h/k=0/1$
in~\eqref{eq:G5-dissection} is one Rademacher arc like every other, and Lemma~\ref{lem:G5-geometry} applies to it with
$k_1=k_2=N$. (These neighbours satisfy $k+k_j=N+1$.)

\begin{lemma}[geometry]\label{lem:G5-geometry}
\begin{itemize}
\item[(b)] $|z'|,|z''|\le\sqrt2k/(N+1)$, and $\Re z',\Re z''\le2k^2/(N+1)^2$.
\item[(c)] On the chord $s_{h,k}=[z',z'']$: $\Re(1/z)\ge1$, and $0<\Re z\le 2k^2/(N+1)^2$. Also
$|s_{h,k}|\le2\sqrt2k/(N+1)$.
\item[(d)] On the arcs of $K$ from $z''$ to $0$ and from $0$ to $z'$ (the complement of the Rademacher arc):
$\Re(1/z)=1$ and $\Re z\le2k^2/(N+1)^2$. Each of these arcs has length at most $(\pi/2)\sqrt2k/(N+1)$.
\end{itemize}
\end{lemma}
\begin{proof}
(b) Farey neighbours satisfy $k+k_j\ge N+1$ (Fact~\ref{fact:neighbours}). So $k^2+k_j^2\ge(N+1)^2/2$. Therefore
$|z'|=k/\sqrt{k^2+k_1^2}\le\sqrt2k/(N+1)$ and $\Re z'=k^2/(k^2+k_1^2)\le2k^2/(N+1)^2$. The same holds for $z''$.

(c) The map $z\mapsto1/z$ sends $K\setminus\{0\}$ onto the line $\Re w=1$, and the open disc onto $\Re w>1$. The chord lies
in the closed disc and avoids $0$. $\Re z$ is affine on the chord and positive at both ends. Finally,
$|s_{h,k}|\le|z'|+|z''|$.

(d) $\Re(1/z)=1$ on $K\setminus\{0\}$. Let $\theta\in(0,\pi]$ be the central angle of the arc from $z'$ to $0$. Its chord
has length $|z'|=\sin(\theta/2)$ and its arc length is $\theta/2\le(\pi/2)\sin(\theta/2)$, by Jordan's inequality. On a
semicircle, $\Re z=(1+\cos\psi)/2$ is monotone, so $\Re z\le\Re z'$ on this arc. The arc from $z''$ is handled the same
way.
\end{proof}

\begin{corollary}\label{cor:G5-exp}
On every chord and every arc in Lemma~\ref{lem:G5-geometry}(d), $|e^{\pi Cz/k^2}|\le e^{2\pi C/(N+1)^2}<e$.
\end{corollary}
\begin{proof}
Lemma~\ref{lem:G5-geometry} and (N2).
\end{proof}

\begin{remark}[the regions of deformation]\label{rem:G5-deform}
Two facts about the region $\Delta_{h,k}$ bounded by the Rademacher arc $z'\to1\to z''$ and the chord $s_{h,k}$ are used
below.
\begin{itemize}
\item \emph{$\Delta_{h,k}$ lies in the half-plane $\Re z\ge\min(\Re z',\Re z'')>0$.} On each of the two semicircular pieces
of the arc, $\Re z$ is monotone, with its maximum at $z=1$. So on the arc, $\Re z\ge\min(\Re z',\Re z'')$. On the chord
this holds by affinity. The region $\Delta_{h,k}$ is a circular segment, which is the convex hull of the arc, so it lies in
that half-plane as well.
\item \emph{The integrands are holomorphic on $\Delta_{h,k}$.} In particular $\Delta_{h,k}$ avoids $z=0$. For $\Re z>0$,
$\tau=h/k+iz/k^2$ has $\Im\tau=\Re z/k^2>0$, so $G_5(e(\tau))$ and the functions derived from it in
\S\ref{sec:G5-transform} are holomorphic there.
\end{itemize}
Cauchy's theorem therefore allows every integral over a Rademacher arc to be replaced by the integral over its chord.
This is used in \S\ref{sec:G5-rem} and \S\ref{sec:G5-ng} alike.
\end{remark}

\subsubsection{The modular transformation at a cusp}\label{sec:G5-transform}

\begin{fact}[Fact~\ref{fact:eta}]\label{fact:G5-eta}
Let $k\ge1$, $(h,k)=1$, $\Re z>0$, and $hH\equiv-1\pmod k$. Put $x=\exp(2\pi ih/k-2\pi z/k^2)$ and
$x'=\exp(2\pi iH/k-2\pi/z)$. Then
\[
F(x)=\omega(h,k)(z/k)^{1/2}\exp\Big(\frac\pi{12z}-\frac{\pi z}{12k^2}\Big)F(x'),
\]
with the principal square root. Equivalently,
$(x;x)_\infty=\omega(h,k)^{-1}(z/k)^{-1/2}\exp(-\frac\pi{12z}+\frac{\pi z}{12k^2})(x';x')_\infty$.
\end{fact}

Let $x=e(\tau)$ with $\tau=h/k+iz/k^2$. Then $G_5(x)=(x;x)_\infty F(x^5)$.

\emph{Case $5\mid k$ (growing cusps), $k=5k_1$.} Write $x^5=\exp(2\pi ih/k_1-2\pi z_1/k_1^2)$ with $z_1=z/5$, and apply
Fact~\ref{fact:G5-eta} to $(h,k_1,z_1)$. Then
\begin{equation}\label{eq:G5-grow}
G_5(x)=\Phi_{h,k}\exp\Big(\frac\pi{3z}-\frac{\pi z}{3k^2}\Big)\hat G_{h,k}(z),
\end{equation}
where $\Phi_{h,k}=\omega(h,k_1)/\omega(h,k)$ is unimodular and $\hat G_{h,k}=(x';x')_\infty F(x_5')$, with
$x_5'=\exp(2\pi iH_1/k_1-10\pi/z)$. The square roots cancel exactly, because $(z_1/k_1)^{1/2}=(z/k)^{1/2}$ is the same
principal root; so no branch choice enters. For the exponents: $-\frac\pi{12z}+\frac{5\pi}{12z}=\frac\pi{3z}$ and
$\frac{\pi z}{12k^2}-\frac{5\pi z}{12k^2}=-\frac{\pi z}{3k^2}$. For $k=5$: $\Phi_{h,5}=w_h$.

\emph{Case $5\nmid k$ (non-growing cusps).} Apply Fact~\ref{fact:G5-eta} to $(5h,k,5z)$ with $5hH_5\equiv-1\pmod k$. Then
\begin{equation}\label{eq:G5-ng}
G_5(x)=\Phi'_{h,k}\sqrt5\exp\Big(-\frac\pi{15z}-\frac{\pi z}{3k^2}\Big)(x';x')_\infty F(x_5''),\qquad
x_5''=\exp(2\pi iH_5/k-2\pi/(5z)),
\end{equation}
with $|\Phi'_{h,k}|=1$. Here $(z/k)^{-1/2}(5z/k)^{1/2}=\sqrt5$, because $z$ and $5z$ have the same argument.

\emph{The shift in $C$.} In~\eqref{eq:G5-dissection}, the factor $e^{-\pi z/(3k^2)}$ combines with $e^{2\pi\nu z/k^2}$ to
give $e^{\pi Cz/k^2}$, where $C=2\nu-1/3$.

Formulae~\eqref{eq:G5-grow} and~\eqref{eq:G5-ng} were checked against direct $q$-series, phases included: $8$ cusps at $60$
digits (relative error $\le4\cdot10^{-60}$), and independently by the referee of this section at $10$ cusps and two values
of $z$ (relative error $\le4.6\cdot10^{-49}$).

\emph{Majorant of $\hat G$.} Put $u=e^{-2\pi/z}$ and $w=\Re(1/z)$. Expanding $(x';x')_\infty=\prod(1-e(mH/k)u^m)$ and
$F(x_5')=\prod(1-e(mH_1/k_1)u^{5m})^{-1}$ gives $\hat G=1+\sum_{j\ge1}g_ju^j$. Coefficientwise,
$|g_j|\le[y^j]\prod(1+y^m)f(y^5)\le[y^j]f(y)f(y^5)=:\phi_j$, because distinct partitions are partitions. Hence
\begin{equation}\label{eq:G5-maj}
|\hat G-1|\le\sum_{j\ge1}\phi_je^{-2\pi jw}.
\end{equation}

\subsubsection{Growing cusps: main term}
By~\eqref{eq:G5-dissection} and~\eqref{eq:G5-grow}, the cusp $h/k$ with $5\mid k$ contributes
$\frac i{k^2}e(-\nu h/k)\Phi_{h,k}[M_{h,k}+R_{h,k}]$, where
\[
M_{h,k}=\int_{z'}^{z''}e^{\pi/(3z)+\pi Cz/k^2}dz,\qquad R_{h,k}=\int_{z'}^{z''}e^{\pi/(3z)+\pi Cz/k^2}(\hat G_{h,k}-1)\,dz.
\]

\begin{fact}[Bessel integral; Fact~\ref{fact:bessel-int}]\label{fact:G5-bessel}
For $x>0$ and $c>0$, $I_1(x)=\frac{x/2}{2\pi i}\int_{c-i\infty}^{c+i\infty}e^{s+x^2/(4s)}s^{-2}\,ds$.
\end{fact}

\emph{From the Hankel loop to the vertical line.} DLMF 10.9.19 \cite{DLMF} gives
$I_1(x)=\frac{x/2}{2\pi i}\int_{-\infty}^{(0+)}e^{s+x^2/(4s)}s^{-2}\,ds$, over a Hankel loop that starts and ends at
$-\infty$ and encircles the negative real axis. For integer order the only singularity of the integrand is the essential
singularity at $s=0$; the loop has no branch cut to respect. Fix $R>c$ and close the region between the loop and the line
$\Re s=c$ with the two arcs $\Gamma_R^\pm$ of $|s|=R$ in $\{\Re s\le c\}$, running from $c\pm i\sqrt{R^2-c^2}$ to the loop.
The region enclosed contains no singularity, so by Cauchy's theorem the loop integral equals the line integral truncated at
height $\sqrt{R^2-c^2}$, plus the integrals over $\Gamma_R^\pm$ (and minus the loop's own parts beyond $|s|=R$). On
$\Gamma_R^\pm$: $|e^s|\le e^c$; $|e^{x^2/(4s)}|\le e^{x^2/(4R)}\le e^{x^2/(4c)}$; $|s^{-2}|=R^{-2}$; the total length is at
most $2\pi R$. So the arcs contribute $O(R^{-1})\to0$. The loop's parts beyond $|s|=R$ vanish too, since $e^s$ decays
along the negative axis. The line integral converges absolutely, because the integrand is $O(|s|^{-2})$ on it. Letting
$R\to\infty$ proves Fact~\ref{fact:G5-bessel}.

\begin{lemma}\label{lem:bessel-contour}
Let $K$ be oriented clockwise, $B>0$, $C>0$. Then $\frac i{k^2}\oint_Ke^{\pi B/z+\pi Cz/k^2}dz=\cP_k$.
\end{lemma}
\begin{proof}
Put $w=1/z=1+i\rho$, $\rho\in\R$. The clockwise traversal of $K\setminus\{0\}$ corresponds to $\rho$ running from $-\infty$
to $\infty$. With $dz=-dw/w^2$, the integral becomes $-\int_{1-i\infty}^{1+i\infty}e^{\pi Bw+\pi C/(k^2w)}w^{-2}dw$.
Substituting $s=\pi Bw$ gives $-\pi B\int e^{s+x^2/(4s)}s^{-2}ds$, with $x=2\pi\sqrt{BC}/k$. By
Fact~\ref{fact:G5-bessel} this equals $-\pi B\cdot 2\pi i\cdot(2/x)I_1(x)$. Multiplying by $i/k^2$ gives
$4\pi^2BI_1(x)/(k^2x)=\cP_k$.
\end{proof}
Numerical check at $k=10$, $\nu=7$: $0.08174002496770177$ on both sides; the referee's check agrees to $30$ digits.
(Lemma~\ref{lem:bessel-contour} is also the identity used for $p=3$ in \S\ref{sec:expansion3}, with $B=\frac12$.)

\emph{Completing the arc.} On $K\setminus\{0\}$ the integrand of $M_{h,k}$ has modulus $e^{\pi/3}|e^{\pi Cz/k^2}|$, which
is bounded, so $\oint_K$ converges absolutely. The complement of the Rademacher arc is the arc $z''\to0\to z'$. By
Lemma~\ref{lem:G5-geometry}(d) and Corollary~\ref{cor:G5-exp}, the integrand there is at most $e^{1+\pi/3}$, and the total
length is at most $\pi\sqrt2k/(N+1)$. Hence
\begin{equation}\label{eq:G5-main}
\frac i{k^2}M_{h,k}=\cP_k(\nu)+\rho_{h,k},\qquad|\rho_{h,k}|\le\frac{\pi\sqrt2e^{1+\pi/3}}{k(N+1)}.
\end{equation}

\emph{Amplitudes.} Put $\fa_k(\nu)=\sum_{(h,k)=1}\Phi_{h,k}e(-\nu h/k)$ for $5\mid k$. This is a sum of $\varphi(k)$
unimodular terms, so $|\fa_k|\le\varphi(k)$, and $\fa_5=\fa$.

\subsubsection{Growing cusps: remainder}\label{sec:G5-rem}
By Remark~\ref{rem:G5-deform}, $R_{h,k}$ may be taken along the chord, where $w\ge1$. By~\eqref{eq:G5-maj},
$|e^{\pi/(3z)}(\hat G-1)|\le\sum_{j\ge1}\phi_je^{(\pi/3-2\pi j)w}$. Each exponent is negative, so each term decreases in
$w$; this is where $B=1/3<2$ is used. The sum is therefore at most its value at $w=1$, namely
$e^{\pi/3}(f(t_0)f(t_0^5)-1)$. With the chord length and Corollary~\ref{cor:G5-exp},
\begin{equation}\label{eq:G5-rem}
\Big|\frac i{k^2}R_{h,k}\Big|\le\frac1{k^2}\cdot\frac{2\sqrt2k}{N+1}e^{1+\pi/3}\big(f(t_0)f(t_0^5)-1\big).
\end{equation}
The $j=1$ term decays, so no correction term needs to be extracted.

\subsubsection{Non-growing cusps}\label{sec:G5-ng}
Here $5\nmid k$. By Remark~\ref{rem:G5-deform} the integral may be moved to the chord. By~\eqref{eq:G5-ng}, for $w\ge1$ on
the chord: $|(x';x')_\infty|\le f(e^{-2\pi w})\le f(e^{-2\pi})$; $|F(x_5'')|\le f(e^{-2\pi w/5})\le f(e^{-2\pi/5})$;
$e^{-\pi w/15}\le e^{-\pi/15}$. Each bound is non-increasing in $w$. With Corollary~\ref{cor:G5-exp}, the integrand is at
most $K_5e$, so each such cusp contributes at most
\begin{equation}\label{eq:G5-ngb}
\frac1{k^2}\cdot\frac{2\sqrt2k}{N+1}\cdot eK_5 .
\end{equation}

\subsubsection{Assembly}
By~\eqref{eq:G5-dissection} and~\eqref{eq:G5-main}--\eqref{eq:G5-ngb},
$s(\nu)-\fa(\nu)\cP_5(\nu)-\sum_{5\mid k,\,10\le k\le N}\fa_k\cP_k=E(\nu)$, where
\[
|E(\nu)|\le\sum_{5\mid k\le N}\frac{\varphi(k)\big[\pi\sqrt2+2\sqrt2(f(t_0)f(t_0^5)-1)\big]e^{1+\pi/3}}{k(N+1)}
+\sum_{5\nmid k\le N}\frac{\varphi(k)\,2\sqrt2eK_5}{k(N+1)}.
\]
For $5\mid k$: $\varphi(k)/k\le4/5$, and there are at most $N/5$ such $k$, so the first weight sum is less than $4/25$. For
$5\nmid k$: $\varphi(k)/k\le1$, and there are at most $N$ values of $k$, so the second weight sum is less than $1$. This
gives $|E|\le E_{\rm arc}+E_{\rm rem}+E_{\rm ng}$. Using $|\fa_k|\le\varphi(k)$ and $N=K(\nu)$ proves
Theorem~\ref{thm:G5}. \qed

\subsubsection{Constants}\label{sec:G5-constants}
Computed at $60$ digits and independently at $50$ digits by the referee:
\begin{center}
\small
\begin{tabular}{lll}
\toprule
constant & value & rounded up\\
\midrule
$K_5$ & $2.85499540499087\ldots$ & $2.8549955$\\
$E_{\rm arc}$ & $5.50644686778894\ldots$ & $5.506447$\\
$E_{\rm rem}$ & $0.00657086323166\ldots$ & $0.006571$\\
$E_{\rm ng}$ & $21.9505238422352\ldots$ & $21.950524$\\
$\Econ$ & $27.4635415732558\ldots$ & $27.4636$\\
\bottomrule
\end{tabular}
\end{center}
The Dedekind sums $s(1,5)=\tfrac15$, $s(2,5)=s(3,5)=0$ and $s(4,5)=-\tfrac15$ were checked at $60$ digits.

\subsubsection{Numerical confirmation and tightness (not used in the proof)}
The referee checked Theorem~\ref{thm:G5} exhaustively over $\nu\in[1,5000]$ and at $1509$ further $\nu\le200000$: no
violation, maximum LHS/RHS $0.9045085$. An earlier ball-arithmetic survey covered every $\nu\le60000$ and every $7$th
$\nu\le200000$ with no failures. The ratio tends to $(5+\sqrt5)/8=0.9045085$ in the class $\nu\equiv2\pmod 5$, where
$|\fa_{10}|=(5+\sqrt5)/2$ is set against $\varphi(10)=4$; so the $\sum\varphi(k)\cP_k$ term must not be weakened in any
later use.

\emph{Uses of Theorem~\ref{thm:G5}.} Write $R(i)$ for its right-hand side. The certificates use
$|s(i)-\fa(i)\cP_5(i)|\le R(i)$ for every $i\ge1$, and exact values of $s(i)$ for $i\le200000$ (\texttt{g5\_coef.py}
$\to$ \texttt{g5\_coeffs.pkl}, cross-checked by an independent recomputation in \texttt{bandfix\_coeffs.py}).

%---------------------------------------------------------------------
\subsection{Exact computation}\label{sec:t-exact}
%---------------------------------------------------------------------

\begin{theorem}[exact]\label{thm:t-exact}
The conjecture holds, for every coefficient and in every class, for all $n\le 979$.
\end{theorem}
\begin{proof}[Method]
The script works modulo $x^{5N_B^2+1}$ and builds $P_{N_A}$ from $P_0=1$. It obtains each next $P_n$ by four exact
shift-subtracts $P\leftarrow P-x^kP$, for $k=5n+1,\dots,5n+4$ with $5\nmid k$. Truncation is exact, because each factor
feeds only higher degrees. It then checks every coefficient with $0\le m\le5n^2$, in every class; a strict wrong sign
counts as a violation, zeros are allowed. The range is run in the segments $1$--$150$, $150$--$291$, $291$--$410$,
$410$--$580$, $580$--$820$, $820$--$920$ and $920$--$1000$, all with $0$ violations; the last run was killed for lack of
memory after the lines $n=920,\dots,979$, all with \texttt{wrong\_sign=0}. The segments abut or overlap, so the exact range
is $1\le n\le979$. Zeros occur only at $m\lesssim5n$, i.e.\ $\mu<2\cdot10^{-3}$, inside the band, where the conjectured
signs are non-strict. Palindromy (Lemma~\ref{lem:palin}) extends each check from $m\le5n^2$ to all $m$.
\cert{centre3\_exact.py N\_A N\_B}{third/logs/exact\_1\_150.log, exact\_150\_291.log, centre3\_exact\_*.log}
\end{proof}

\emph{Further exact checks (used in \S\ref{sec:t-band}).}
\begin{itemize}
\item Every class and every $m<4N$, for every $n\in[291,2000]$. The method is $c_n=G_5E_n\bmod q^{4N}$ with
$6E_n=6+6p_1+3(p_1^2+p_1(q^2))+(p_1^3+3p_1(q^2)p_1+2p_1(q^3))$ (the cycle index of $S_2$ and $S_3$), where
$p_1=\sum_{N\le k<4N,\,5\nmid k}q^k$, computed with \texttt{fmpz\_poly.mul\_low}. The script was validated against the full
product for $n\le14$, and it reports zero sign failures. An audit recomputed $16$ values of $n$ in $[291,2000]$ by a direct
product that does not use $E_n$ or the coefficient table, with identical results for all $m<4N$.
\cert{bandfix\_exact.py 291 2000}{third/logs/bandfix\_exact.log}
\item The coefficients $s(i)$ for $i\le200000$ come from \texttt{g5\_coef.py} (pentagonal series times $p(k)$ at $q^5$).
They were recomputed independently.
\cert{bandfix\_coeffs.py}{third/logs/bandfix\_coeffs.log}
\end{itemize}

\begin{remark}
The analytic certificates below hold from $n\ge821$ (Laplace classes $3$--$4$ and the centre) or from $n\ge291$ (Laplace
classes $0$--$2$). The overlap with the exact range is therefore $821\le n\le979$. From here on, $n\ge980$.
\end{remark}

%---------------------------------------------------------------------
\subsection{The band, \texorpdfstring{$\lambda<4$}{lambda < 4}}\label{sec:t-band}
%---------------------------------------------------------------------

Every part of $E_n$ is $\ge N$. So for $m<4N$ at most three parts contribute, and $c(m)=s(m)+b_1+b_2+b_3$, where $b_j$
collects the $j$-part terms. For $j\in[N,2N)$, $e(j)=1$ if $5\nmid j$ and $e(j)=0$ otherwise.

\subsubsection{Classes 3 and 4}\label{sec:t-band34}

\emph{One-part block.} Let $m\equiv3$ or $4\pmod 5$ and write $m=5M+r$. Then $s(m)=0$, and
$b_1=\sum_{j\in[N,m],\,5\nmid j}s(m-j)=\sum_{i\le m-N,\ i\not\equiv m}s(i)$. Since $s(i)=0$ for $i\equiv3,4$, only
$i\equiv0,1,2$ survive. Also $m-N=5(M-n)+r-1$ with $r-1\in\{2,3\}$, so every block $\{5t,5t+1,5t+2\}$ with $t\le T:=M-n$
is included in full. Hence
\[
b_1=G(T)=\sum_{t\le T}g(t),\qquad g(t)=s(5t)+s(5t+1)+s(5t+2).
\]

\begin{lemma}\label{lem:t-L1}
The main term satisfies $\mathrm{main}(g(t))\le-5\cP_5'(5t)$. Moreover $g(t)<-4\cP_5'(5t)<0$ for $t\ge200$.
\end{lemma}
\begin{proof}
The main term of $g(t)$ is $\fa(0)\cP_5(5t)+\fa(1)\cP_5(5t+1)+\fa(2)\cP_5(5t+2)$. Since $\fa(0)=|\fa(1)|+|\fa(2)|$, this
equals $-|\fa(1)|\,[\cP_5(5t+1)-\cP_5(5t)]-|\fa(2)|\,[\cP_5(5t+2)-\cP_5(5t)]$. By the series of $I_1$,
$\cP_5(i)=\sum_{r\ge1}a_5^r x^{r-1}/(r!\,(r-1)!)$ is a power series with positive coefficients in $x=i-\tfrac16$ (not in
$i$ itself). So $\cP_5'$ is increasing for $x>0$, and the main term is at most
$-(|\fa(1)|+2|\fa(2)|)\cP_5'(5t)=-5\cP_5'(5t)$.

By Theorem~\ref{thm:G5}, $|g(t)-\mathrm{main}|\le3R_{\rm grp}(5t+2)$, where $R_{\rm grp}$ is the (increasing) right-hand
side of Theorem~\ref{thm:G5}. Lemma~L1 states $3R_{\rm grp}(X+2)\le0.018\,\cP_5'(X)$ for all $X\ge1000$. Its proof has four
steps:
\begin{enumerate}
\item write $\cP_5'(X)=4a_5^2I_2(Y)/Y^2$ with $Y=2\sqrt{a_5(X-1/6)}$;
\item use $I_2\ge I_{5/2}$ (monotonicity in the order, Fact~\ref{fact:bessel}) and the closed form of $I_{5/2}$, which give
$I_2(Y)\ge e^Y(1-3/Y-2e^{-2Y})/\sqrt{2\pi Y}$;
\item use $I_1\le I_{1/2}$, $\cP_k\le\cP_{10}$ for $k\ge10$, and $\sum_{5\mid k\le K}\varphi(k)\le K^2/10+K/2$;
\item bound the ratio by $c_1X^{3/2}e^{-Y/2}+c_2Y^{5/2}e^{-Y}$. Both terms decrease for $X\ge1000$ (sign of the
log-derivative), so the bound evaluated at $X=1000$ in Arb holds for all larger $X$: the log prints
\texttt{L1: sup\_\{X>=1000\} 3 Rgrp(X+2)/P5'(X) <= [0.017972 +/- 1.12e-7]}.
\end{enumerate}
This gives $g(t)\le-5\cP_5'+0.018\cP_5'<-4\cP_5'$ for $5t\ge1000$.
\cert{bandfix\_analytic.py 10001}{third/logs/bandfix\_analytic.log}
\end{proof}

\begin{lemma}[exact]\label{lem:t-block-exact}
For $t\le39999$, $g(t)<0$ except at $g(1)=1$ and $g(5)=0$. Every partial sum satisfies $G(T)\le0$, with equality only at
$T=1$. The computation uses the table of exact $s(i)$, which the same script recomputes independently first.
\cert{bandfix\_coeffs.py}{third/logs/bandfix\_coeffs.log}
\end{lemma}

\begin{corollary}\label{cor:t-block}
$G(T)\le0$ for every $T$, and $G(T)\le g(T)<0$ for $T\ge200$. Hence, for classes $3$ and $4$ and every $n$: for $m<N$,
$c(m)=s(m)=0$; for $N\le m<2N$, $c(m)=G(T)\le0$.
\end{corollary}

\emph{The range $2N\le m<4N$.}
\begin{itemize}
\item \emph{For $980\le n\le2000$ (exact):} by the exact band run, which covers $291\le n\le2000$. Over that whole run the
log prints the largest ratio $|b_2+b_3|/|b_1|$ as \texttt{0.0001}, at $n=291$, $m=5819$; an audit's rerun gives the
unrounded value $1.1302\cdot10^{-4}$.
\item \emph{For $n>2000$ (so $N\ge10001$).} Write $I=m-2N$. Then $5T\ge I+N-3$.
\begin{itemize}
\item \emph{Denominator.} $|b_1|\ge|g(T)|\ge4\cP_5'(I+N-3)$, by Corollary~\ref{cor:t-block} and the monotonicity of
$\cP_5'$.
\item \emph{Numerator.} $|b_2|+|b_3|\le\sum_{d\le I}w(2N+d)|s(I-d)|$. Here the number of pairs is at most $d/2+1$, and the
number of triples at $3N+e$ is at most $(e+3)^2/12+1/2$. Also $|s(i)|\le e^{2\sqrt{a_5i}}$: this is exact for $i\le200000$,
and beyond that Theorem~\ref{thm:G5} gives the stronger constant $7.7\cdot10^{-5}$.
\item \emph{Split.} Split the sum at $D=\lfloor3\sqrt{2N/a_5}\log N\rfloor$. The exponent
$2\sqrt{a_5}(\sqrt I-\sqrt{I+N-19/6})$ increases in $I$, so it is largest at $I=2N$. This gives
\[
\frac{|b_2|+|b_3|}{|b_1|}\le\Psi(N)=C_\Psi\,N^{9/4}(\log N)^2e^{-2\sqrt{a_5}(\sqrt3-\sqrt2)\sqrt N}.
\]
\item \emph{Conclusion.} $d\log\Psi/dN<0$ exactly when $(c_c/2)\sqrt N>9/4+2/\log N$, where
$c_c=2\sqrt{a_5}(\sqrt3-\sqrt2)$ and $c_c/2=0.16304$. This first holds at $N=257$ and persists for larger $N$, so $\Psi$ is
decreasing for $N\ge257$. In Arb, $\Psi(10001)\le0.19769$. Hence $c(m)\le b_1(1-0.1977)<0$.
\end{itemize}
\cert{bandfix\_analytic.py 10001}{third/logs/bandfix\_analytic.log}
\end{itemize}

\begin{corollary}\label{cor:t-band34}
For classes $3$ and $4$, every $n\ge980$ and every $m<4N$: $c(m)\le0$.
\end{corollary}

\subsubsection{Classes 0, 1, 2}\label{sec:t-band012}
For $980\le n\le2000$ the band is covered by the exact run. For $n>2000$ write $c(m)=s(m)+\sum_{j\ge N}e(j)s(m-j)$. The
weights satisfy $e(j)\le1$ for $j<2N$ and $e(j)\le1+(N+1)+((N+3)^2/12+1/2)\le N^2\le m^2$ for $j<4N$. With
$A(I)=\sum_{i\le I}|s(i)|$, this gives
\[
|c(m)-s(m)|\le A(m-N)+m^2A(m-2N)\le A(m-N_{\min})+m^2A(m-2N_{\min}),
\]
where $N_{\min}(m)=\max(10001,\lfloor m/4\rfloor+1)$.
The second inequality uses that $A$ is non-decreasing and that $m<4N$ is equivalent to $N\ge\lfloor m/4\rfloor+1$. So one
inequality per $m$ covers every $n>2000$ at once.

\begin{lemma}\label{lem:t-band012}
\begin{enumerate}
\item[(A)] (\emph{exact}, $m\le200000$.) $\sgn s(m)=\sgn\fa(m)$ and $|s(m)|>A(m-N_{\min})+m^2A(m-2N_{\min})$ for every $m$
in classes $0,1,2$, except $s(7)=0$. That coefficient has $m<N$, so $c(7)=0$, which is allowed. The worst ratio is
$5.06\cdot10^{-10}$.
\item[(B)] (\emph{certified}, $m>200000$.) With the inputs $|s(i)|\le e^{c\sqrt i}$ ($c=2\sqrt{a_5}$),
$\sum_{i\le I}e^{c\sqrt i}\le e^{c\sqrt I}(1+2\sqrt I/c)$, $\cP_5(m)\ge h_0m^{-3/4}e^{c\sqrt m}$ (via $I_1\ge I_{3/2}$) and
$R(m)\le1.3\,m\,e^{c\sqrt m/2}+\Econ$, one has
$[\text{corrections}+R(m)]/((5-\sqrt5)/2\cdot \cP_5(m))\le Q\le5.14\cdot10^{-20}$. Each term has the form
$C'm^{p'}e^{-\kappa\sqrt m}$ with $p'\in\{1.25,3.25,1.75,0.75\}$, and is decreasing once $\sqrt m>2p'/\kappa$.
\end{enumerate}
\cert{cls012\_band.py}{third/logs/cls012\_band.log}
\end{lemma}

\begin{corollary}\label{cor:t-band012}
For classes $0,1,2$, every $n\ge980$ and every $m<4N$, $c(m)$ has the conjectured sign.
\end{corollary}

%---------------------------------------------------------------------
\subsection{The Laplace region, \texorpdfstring{$\lambda\ge4$, $\mu<0.105$}{lambda >= 4, mu < 0.105}}\label{sec:t-laplace}
%---------------------------------------------------------------------

Both certificates of this subsection share one exact representation (\S\ref{sec:t-repr}) and one Euler--Maclaurin lemma
(\S\ref{sec:t-em}). They differ in the saddle, because classes $3$ and $4$ are thin. Write $\vp_\eta=1+i\eta$.

\subsubsection{Exact Bessel--Laplace representation}\label{sec:t-repr}
Write $\fa(i)=\sum_{h=1}^4A_h\omega_5^{hi}$ with $A_1=e^{i\pi/5}$, $A_4=e^{-i\pi/5}$ and $A_2=A_3=1$. By
Lemma~\ref{lem:factor} and Theorem~\ref{thm:G5},
\[
c(m)=e(m)+\sum_hA_h\omega_5^{hm}J_h+\rho,\qquad J_h=\sum_{j<m}e(j)\omega_5^{-hj}\cP_5(m-j),\qquad|\rho|\le\sum_{j<m}e(j)R(m-j).
\]
Here $e(m)=e(m)s(0)$ is the $j=m$ term, and $\omega_5^{-hj}\omega_5^{hm}=\omega_5^{h(m-j)}$ reproduces $\fa(m-j)$.

For $x=i-1/6>0$ we have $\cP_5(i)=\sqrt{a_5/x}\,I_1(2\sqrt{a_5x})$, which is the inverse Laplace transform of
$e^{a_5/u}-1$. Put $r_1(u)=e^{a_5/u}-1-a_5/u=O(|u|^{-2})$. Then
$\frac1{2\pi i}\int_{(W)}r_1(u)e^{xu}du=\cP_5-a_5$ for $x>0$ and $=0$ for $x<0$. Since $r_1=O(|u|^{-2})$, Fubini applies,
and for every $W>0$, exactly,
\[
J_h=a_5H_h+K_h,\qquad H_h=\sum_{j\le m-1}e(j)\omega_5^{-hj},
\]
\[
K_h=\frac1{2\pi}\int_{\R}r_1(W+iy)\,e^{(m-1/6)(W+iy)}E_n(\omega_5^{-h}e^{-W-iy})\,dy.
\]

\subsubsection{Euler--Maclaurin lemma for the twisted \texorpdfstring{$E_n$}{En}}\label{sec:t-em}

\begin{lemma}\label{lem:t-em}
Let $\theta\in(0,1)$, and let $f$ be such that $f,f'\to0$ at $\infty$ and $f''$ is integrable. Then
\[
\sum_{t\ge n}f(t+\theta)=\int_n^\infty f-B_1(\theta)f(n)+R,\qquad|R|\le\tfrac1{12}\Big(|f'(n)|+\int_n^\infty|f''|\Big).
\]
\end{lemma}
\begin{proof}
On each unit interval $[t,t+1]$, write $f(t+\theta)-\int_t^{t+1}f=\int_t^{t+1}k(x-t)f'(x)\,dx$, where $k(x)=x-H(x-\theta)$
and $H$ is the Heaviside function. The mean of $k$ is $B_1(\theta)$, and $k-B_1(\theta)=\tilde B_1(x-\theta)$. The mean
part telescopes to $-B_1(\theta)f(n)$. Integrate the remaining part by parts once more, with kernel
$P(x)=(\tilde B_2(x-\theta)-B_2(\theta))/2$. Its constant part contributes $(B_2(\theta)/2)f'(n)$, with
$|B_2(c/5)|/2\le 11/300<1/12$. Its periodic part is bounded by $\sup|\tilde B_2|/2=1/12$ times $\int|f''|$.
\end{proof}
The constant part must be split off: without the split, $\sup|P|$ reaches $0.12>1/12$ at $\theta=2/5$.

\emph{Application.} Write $k=5t+c$ with $c\in\{1,2,3,4\}$. Put $\Phi_c(s)=\log(1-\omega_5^{-hc}e^{-s})$, $u=W\vp_\eta$,
$\sigma=5nu$, $V=5nW$ and $\xi=e^{-V}$. Apply Lemma~\ref{lem:t-em} to $f(t)=\Phi_c(5ut)$ with $\theta=c/5$, and sum over
$c$. This gives
\[
\log E_n(\omega_5^{-h}e^{-u})=-\gamma(\sigma)/u+\Omega_h(\sigma)+R_h,
\]
where $\gamma(z)=\big[\Li_2(e^{-z})-\Li_2(e^{-5z})/5\big]/5$;
$Q(s)=\sum_c\Phi_c(s)=\log(1-e^{-5s})-\log(1-e^{-s})$, independent of $h$ because
$\prod_c(1-\omega_5^{-hc}w)=(1-w^5)/(1-w)$; $\Omega_h(\sigma)=\sum_cB_1(c/5)\Phi_c(\sigma)$; and
\[
|R_h|\le\bar R:=(5W/3)|\vp_\eta|\xi\big[1/(1-\xi)+|\vp_\eta|/(1-\xi)^2\big].
\]
\emph{Derivation of $\bar R$.} Here $f(t)=\Phi_c(5ut)$, so $f'=5u\,\Phi_c'$ and $f''=25u^2\Phi_c''$, with
$|u|=W|\vp_\eta|$ and $\Re(5ut)=5Wt$. From $\Phi_c'(s)=\omega_5^{-hc}e^{-s}/(1-\omega_5^{-hc}e^{-s})$,
$\Phi_c''(s)=-\omega_5^{-hc}e^{-s}/(1-\omega_5^{-hc}e^{-s})^2$ and $|1-\omega_5^{-hc}e^{-s}|\ge1-e^{-\Re s}$ we get
$|\Phi_c'(s)|\le e^{-\Re s}/(1-e^{-\Re s})$, so $|f'(n)|\le5W|\vp_\eta|\,\xi/(1-\xi)$ since $\Re(5un)=V$; and
$|\Phi_c''(s)|\le e^{-\Re s}/(1-e^{-\Re s})^2$, so, as the denominator is at least $(1-\xi)^2$ for $t\ge n$,
\[
\int_n^\infty|f''|\,dt\le\frac{25W^2|\vp_\eta|^2}{(1-\xi)^2}\int_n^\infty e^{-5Wt}dt=\frac{5W|\vp_\eta|^2\xi}{(1-\xi)^2}.
\]
Lemma~\ref{lem:t-em} gives $\tfrac1{12}(|f'(n)|+\int|f''|)$ for each $c$; summing over the four values of $c$ gives $\bar R$.

The pairs $c\leftrightarrow5-c$ have opposite $B_1$ and conjugate roots of unity, so $\Re\Omega_h(W)=0$, i.e.\
$|e^{\Omega_h(W)}|=1$. Also $|\Omega_h|\le\bar\omega:=0.8\,\Li_1(\xi)$, since $\sum_c|B_1(c/5)|=0.8$. (In
\texttt{cls012\_em\_check.py} and \texttt{lap34\_emchk.py} the bound exceeds the true remainder by a factor $5$--$20$.)

\subsubsection{Classes 0, 1, 2}\label{sec:lap012}

\emph{Saddle.} Define $\alpha(V)=a_5-\gamma(V)-VQ(V)/5$ and choose $W$ by $(m-1/6)W^2=\alpha(V)$, where $V=5nW$. With
$X=1/W$ the exponent becomes
\[
\frac{a_5}u+(m-\tfrac16)u-\frac{\gamma(\sigma)}u=X\,\Xi_V(\eta),\qquad\Xi_V(\eta)=\frac{a_5-\gamma(V\vp_\eta)}{\vp_\eta}+\alpha(V)\vp_\eta .
\]
$\Xi_V$ depends on $V$ only. It satisfies $\Xi_V'(0)=0$ identically (this is the definition of $\alpha$), and
$\Xi(-\eta)=\overline{\Xi(\eta)}$. Put $\Xi_2=-\Xi''(0)>0$ and $\Amp=W^{3/2}e^{X\Xi(0)}/\sqrt{2\pi\Xi_2}$. Then
\[
c(m)=\Amp\Big[M+\sum_hA_h\omega_5^{hm}e^{\Omega_h}\delta_h\Big]+E_2',\qquad
M=\sum_hA_h\omega_5^{h\cdot\mathrm{cls}}e^{\Omega_h(W)}\in\R .
\]
As $V\to\infty$, $M\to\fa(\mathrm{cls})$. $M$ is the main term with the correct normalisation $|E_n(\omega_5e^{-W})|$.

\emph{Error terms.} All are relative to $\Amp$, and each is either independent of $X$ or non-increasing in $X$. Set
$\eta_0=c_c'/\sqrt{X\Xi_2}$ with $c_c'=4$, and $s=\eta\sqrt{X\Xi_2}$.
\begin{itemize}
\item \emph{Window, $|\eta|\le\eta_0$.} Write $\psi=X(ic_3\eta^3+R_4)+\Delta\Omega+R_h$. The cubic term is a pure phase, so
$|e^\psi-1|\le\min(k_3s^3,2)e^r+e^r-1$ with $r=k_4s^4+k_1s+\bar R$, where $k_3=|c_3|/(\Xi_2^{3/2}\sqrt X)$,
$k_4=\sup|\Xi''''|/(24\Xi_2^2X)$ and $k_1=\beta_1/\sqrt{X\Xi_2}$. Here $\beta_1=0.8V\xi/(1-\xi)$ bounds $|d\Omega/d\eta|$.
The supremum of $|\Xi''''|$ is enclosed on sub-balls of width $0.01$ covering $[0,\eta_0(X_{\min})]$. The Gaussian
truncation contributes $\erfc(c_c'/\sqrt2)$.
\item \emph{Mid region, $\eta_0<|\eta|\le2$.} $D(\eta)=\Re(\Xi(0)-\Xi(\eta))\ge d_k\eta^2$ on $170$ fixed $\eta$-tiles,
using the better of the Taylor bound $\Xi_2/2-\sup|\Xi''''|\eta^2/24$ and a direct ball evaluation. The contribution is
$e^{2\bar\omega+\bar R}\sum_k\sqrt{\Xi_2/(2d_k)}\,\erfc\big(\max(c_c'\sqrt{d_k/\Xi_2},\ \eta_k\sqrt{X_{\min}d_k})\big)$.
\item \emph{Crude terms $E_2$.} These are: the $(1+a_5/u)$ part of $r_1$; $|y|\in[2W,1]$, with
$|r_1|\le e^{a_5X/(1+\eta_1^2)}+1+a_5X/\eta_1$; $|y|>1$, with $|r_1|\le(a_5^2/2y^2)e^{a_5}$; $e(m)\le e^{mW}E_n(e^{-W})$;
$a_5|H_h|\le a_5e^{(m-1)W}E_n(e^{-W})$; and the $\rho$-term, by Rankin's trick at the same $W$, with
$\cP_k(i)\le W(e^{a_kX}-1)e^{(i-1/6)W}$ and $k\le K(m)$.
\end{itemize}

\emph{The $\rho$-term in detail.} By \S\ref{sec:t-repr}, $|\rho|\le\sum_{j<m}e(j)R(m-j)$, where
$R(i)=\sum_{5\mid k,\,10\le k\le K(i)}\varphi(k)\cP_k(i)+\Econ$. Every $i=m-j$ here is $\ge1$, so Theorem~\ref{thm:G5}
applies. Write $x=i-\tfrac16$ and $a_k=2\pi^2/(3k^2)$, so that $a_{10}=a_5/4$ and $a_k\le a_5/9$ for $k\ge15$.
\begin{itemize}
\item \emph{Bessel part.} The series of $I_1$ gives $\cP_k(i)=\sum_{r\ge1}a_k^rx^{r-1}/(r!(r-1)!)$. With
$x^{r-1}/(r-1)!\le e^{xW}/W^{r-1}$ this gives $\cP_k(i)\le W(e^{a_kX}-1)e^{(i-1/6)W}$.
\item \emph{Counting.} $\varphi(10)=4$, and $\sum_{j=3}^{K/5}\varphi(5j)\le\sum4j\le2\bar K(\bar K+1)$, where
$\bar K=K/5$ and $K$ bounds $K(i)$ for every $i\le m$. The code uses $K=\sqrt{4\pi(a_5X^2+1)}+2$, since
$m-\tfrac16=\alpha X^2\le a_5X^2$.
\item \emph{Rankin.} Since $e(j)\ge0$, $\sum_{j<m}e(j)e^{(m-j-1/6)W}\le e^{(m-1/6)W}E_n(e^{-W})$. So the Bessel part
contributes at most $W\big[4e^{a_5X/4}+2\bar K(\bar K+1)e^{a_5X/9}\big]e^{(m-1/6)W}E_n(e^{-W})$.
\item \emph{$\Econ$ part.} $e^{(m-j)W}\ge1$ for $j\le m$, so
$\sum_{j<m}e(j)\Econ\le \Econ e^{mW}E_n(e^{-W})=\Econ e^{W/6}\cdot e^{(m-1/6)W}E_n(e^{-W})$.
\end{itemize}
Together these give the term
$T_\rho=F_m e^{W/6}\big(W[4e^{a_5X/4}+e^{a_5X/9}2\bar K(\bar K+1)]+\Econ\big)$, where $F_m$ is an upper bound for
$e^{(m-1/6)W}E_n(e^{-W})/\Amp$. The script also applies the factor $e^{W/6}$ to the Bessel part, where it is not needed;
this only enlarges the bound.

All the crude terms use $E_n(e^{-W})\le\exp(0.8X\Li_2(\xi)+4\Li_1(\xi))$, which holds because $E_n$ has non-negative
coefficients, so $|E_n(\omega_5^{-h}e^{-u})|\le E_n(e^{-W})$. The total is at most $\sqrt{2\pi\Xi_2}X^{3/2}e^{-cX}\mathrm{poly}(X)$,
with $c\ge\kappa-a_5/4$ and $\kappa=a_5-\gamma-0.8\Li_2(\xi)$; the rate $a_5/4$ comes from the $k=10$ Rankin term. The
script asserts $c>0$ and $X_{\min}>4.5/c$, and $E_2\le2.3\cdot10^{-10}$.

\emph{Certified condition.} $\sgn(\fa)M-4\epsilon_{\rm tot}-E_2>0$, where
$\epsilon_{\rm tot}=\varepsilon_{\rm win}+\erfc(c_c'/\sqrt2)+\varepsilon_{\rm mid}$ (called \texttt{delta} in the log).
Each factor $A_h\omega_5^{hm}e^{\Omega_h(W)}$ is unimodular, because $\Re\Omega_h(W)=0$, and $|\delta_h|\le\epsilon_{\rm tot}$.
So the bracket is $M$ plus at most $4\epsilon_{\rm tot}$, and the condition implies that $c(m)$ has the sign of
$\fa(\mathrm{cls})$.

\emph{Range of $X$.} $X=5n/V\ge1455/V$ for $n\ge291$. $\lambda\ge4$ gives $\alpha X^2\ge4VX+23/6$, so $X\ge4V/a_5$. Per
tile, $X_{\min}=\max(1455/V_{hi},4V_{lo}/a_5)$.

\emph{Monotonicity in $X$.} $k_3,k_1\propto X^{-1/2}$ and $k_4\propto X^{-1}$; $\bar R\propto W$; the erfc arguments are
non-decreasing in $X$; the window sets are taken at $\eta_0(X_{\min})$; the crude terms decay. So one evaluation at
$X_{\min}$ per tile covers every $X\ge X_{\min}$, that is, every $(n,m)$ whose saddle lies in the tile. No monotonicity in
$\lambda$ is used.

\begin{lemma}[classes $0$--$2$, tiles]\label{lem:lap012-tiles}
On $500$ contiguous tiles of width $0.02$ covering $V\in[2,12]$ (the last tile snapped to $[11.98,12]$), the certified
condition holds, with $0$ failures. The minimum margins $\sgn(\fa)M-4\epsilon_{\rm tot}-E_2$ are $3.0700$ (class $0$,
$|\fa|=3.618$), $1.6879$ (class $1$, $|\fa|=2.236$) and $0.8339$ (class $2$, $|\fa|=1.382$). The worst $\epsilon_{\rm tot}$ is
$0.137$, at $V=9.78$, $X_{\min}=148.6$.
\cert{cls012\_lap.py 2.0 12.0 0.02}{third/logs/cls012\_lap\_rerun.log (line 1)}
\end{lemma}

\emph{Tail cell $V\ge11.99$.} Write $\Xi=a_5/\vp_\eta+a_5\vp_\eta+\Delta$. On the strip $|\Im\eta|\le1/2$ we have
$\Re\vp_\eta\ge1/2$, so $|\gamma(V\vp_\eta)/\vp_\eta|\le\frac{12}{25}\Li_2(e^{-V/2})$. Cauchy estimates on discs of radius
$1/2$ give $\Xi_2\in2a_5\pm8\varepsilon$, $|c_3|\le a_5+8\varepsilon$ and $|\Xi''''|\le24a_5+384\varepsilon$. The
mid-region bound is $a_5\eta^2/(1+\eta^2)-\frac{12}{25}\Li_2$. Also $M\ge|\fa|-4(e^{\bar\omega}-1)$ and
$X\ge4V/a_5\ge4V_T/a_5$. Every $V$-dependent quantity moves in the safe direction as $V$ grows, so a single evaluation at
$V_T$ covers all $V\ge V_T$.

\begin{lemma}[classes $0$--$2$, tail cell]\label{lem:lap012-tail}
For all $V\ge V_T=11.99$ the certified condition holds, with margins $3.0950$, $1.7130$, $0.8589$ (classes $0,1,2$) and
$X_{\min}=182.2$. The cell overlaps the last tile on $[11.99,12]$, so coverage at $V=12$ does not depend on the snap.
\cert{cls012\_lap.py tail 11.99}{third/logs/cls012\_lap\_rerun.log (line 2)}
\end{lemma}

\emph{Validation (not part of the proof).} \texttt{cls012\_asym\_check.py}: $c/(\Amp\cdot M)=0.990$--$0.998$ at
$n=60,100$. An audit checked $45$ exact points, $n\in\{291,437,500,700,1500,2500\}$, $V=2.08$--$40$, including
$\lambda\approx4$ and $\mu\approx0.105$: every sign correct; $|c/\Amp-M|$ ranged over $0.0005$--$0.0185$, against a
certified bound of $0.27$--$0.54$.

\subsubsection{Classes 3 and 4: the re-centred (thin) saddle}\label{sec:lap34}

For classes $3$ and $4$, $\sum_hA_h\omega_5^{h\cdot\mathrm{cls}}=\fa(\mathrm{cls})=0$, so the constant term of the twisted sum
cancels. The survivor carries $e^{-5nu}=e^{-\sigma}$, which is an $O(N)$ term \emph{in the exponent}. It must not be
treated as an amplitude (see Appendix~\ref{app:withdrawn}); here it is moved into the exponent.

\emph{The thin function.} Put
\[
g(s)=e^{s}\sum_{h=1}^4A_h\omega_5^{h\cdot\mathrm{cls}}\exp(\Omega_h(s)),\qquad s=\sigma=V\vp_\eta .
\]
Since $\omega_5^{hm}=\omega_5^{h\cdot\mathrm{cls}}$, we have $\sum_hA_h\omega_5^{hm}e^{\Omega_h(\sigma)}=e^{-\sigma}g(\sigma)$ exactly.

\begin{lemma}[the thin function]\label{lem:thin}
\begin{enumerate}
\item \emph{Leading coefficient.} Expand $e^{\Omega_h}=1+\Omega_h+\dots$. The constant term cancels. The coefficient of
$e^{-s}$ is $f_\infty=-\sum_hA_h\omega_5^{h\cdot\mathrm{cls}}\sum_cB_1(c/5)\omega_5^{-hc}=-1$ exactly, for both classes
(checked in Arb to $10^{-37}$, \texttt{lap34\_lib.finf}).
\item \emph{Uniform bound.} For every $s$ with $\Re s=V$,
$|g(s)-f_\infty|\le\varepsilon_g(V):=4e^V\big[0.8(\Li_1(\xi)-\xi)+\bar\omega^2e^{\bar\omega}/2\big]$. This follows from
$|\Omega_h-\Omega_h^{(1)}|\le0.8\sum_{r\ge2}\xi^r/r$ and $|e^\Omega-1-\Omega|\le|\Omega|^2e^{|\Omega|}/2$, where
$\Omega_h^{(1)}$ is the first-order term. $\varepsilon_g$ is $\xi^{-1}$ times a series in $\xi$ with non-negative
coefficients starting at $\xi^2$, so it decreases in $V$. For example $\varepsilon_g(11.99)=1.787\cdot10^{-5}$.
\item \emph{Reality.} $g(V)$ is real for real $V$, because $h\leftrightarrow5-h$ are conjugate. The Arb enclosures of $g$ on
the $V$-tiles lie in $[-1.13,-0.96]$.
\end{enumerate}
\end{lemma}
(An audit tested (2) at $6000$ points and found $\max|g+1|/\varepsilon_g=0.233$.)

\emph{Saddle.} Choose $W$ by $\alpha(V)X^2=m-\tfrac16-5n$. The exponent $a_5/u-\gamma/u+(m-1/6)u-5nu$ then equals
$X\Xi_V(\eta)$, with the same $\Xi_V$ as in \S\ref{sec:lap012}. So
\[
c(m)=\Amp\,[g(V)+\Delta],\qquad \Amp=W^{3/2}e^{X\Xi_V(0)}/\sqrt{2\pi\Xi_2}>0,
\]
and the certificate proves $|\Delta|<-g(V)$, which gives $c(m)<0$. All the $O(N)$ parts of $\log F_m$ are inside $X\Xi_V$:
the kernel, the twisted-product term $-\gamma/u$, and the thin shift $-5nu$. The remaining $\log g(V\vp_\eta)$ is $O(1)$.
It is kept out of the Gaussian and paid for in the window error through $\sup|g'|$.

\emph{Error terms} (relative to $\Amp$, all non-increasing in $X$).
\begin{itemize}
\item \emph{Window.} With $G(\eta)=e^{V\vp_\eta}\sum_hA_h\omega_5^{hm}e^{\Omega_h+R_h}$, $\psi=X(c_3\eta^3+R_4)$ and $c_3$
purely imaginary,
\[
|e^\psi G-g(V)|\le e^{k_4s^4}(V|\eta|\sup_{\rm win}|g'|+\varepsilon_R)+|g(V)|\big(\min(k_3s^3,2)e^{k_4s^4}+e^{k_4s^4}-1\big).
\]
Here $\varepsilon_R=4e^{\bar\omega}e^V(e^{\bar R}-1)\ge|G-g(V\vp_\eta)|$; $\sup|g'|$ is the smaller of a ball evaluation
(sub-balls of width at most $0.05$ in $\Im s$) and the Cauchy bound $\varepsilon_g(V-1)$; the integral is an upper Riemann
sum over $800$ $s$-cells; the Gaussian truncation contributes $|g|\erfc(c_c'/\sqrt2)$.
\item \emph{Mid region, $\eta_0<|\eta|\le2$.} The bound $\Re(\Xi(0)-\Xi(\eta))\ge d_k\eta^2$ on $170$ $\eta$-tiles. On each
tile, $\sup|G|\le\sup|g(V\vp_\eta)|+\varepsilon_R$. Tiles whose erfc argument exceeds $40$ are bounded jointly, via
$4e^{V+\bar\omega+\bar R}\cdot170\sqrt{\Xi_2/2d_{\min}}\erfc(40)$.
\item \emph{Crude terms.} As in \S\ref{sec:lap012}, but with an extra factor $e^V$, because $\Amp$ carries $e^{-V}$. The rate
is $\kappa=a_5-\gamma-0.8\Li_2(\xi)$. Tiles use $e^{V_{hi}}$; the tail cell uses $V\le a_5X/3$. The script asserts decay
$c\ge\kappa-\max(a_5/5,a_5/4)$ (tiles), or $c\ge\kappa-a_5/3-a_5/4$ (tail), and $X_{\min}>4.5/c$. $E_2\le10^{-8}$
everywhere.
\end{itemize}

\emph{Range of $X$.} $X=5n/V\ge4105/V$ for $n\ge821$. $\lambda\ge4$ gives $m-1/6-5n\ge15n+23/6>3VX$, so $X>3V/a_5$. Per
tile, $X_{\min}=\max(4105/V_{hi},3V_{lo}/a_5)$. For the tail cell, $X\ge X^*=\sqrt{15\cdot821/a_5}=216.3$.

\begin{lemma}[classes $3$--$4$]\label{lem:lap34}
\begin{enumerate}
\item On $500$ tiles of width $0.02$ covering $V\in[2,12]$ (last tile snapped to $[11.98,12]$), $|\Delta|<-g(V)$ with $0$
failures; the minimum margins $-g-|\Delta|$ are $0.8363$ (class $3$) and $0.8361$ (class $4$), on the tile $[11.98,12]$
(window $0.119$, mid $0.005$, $\varepsilon_R$ $0.042$).
\item On the analytic tail cell $V\ge V_T=11.99$, the margin is $0.8146$ for both classes ($0.81456$). The tail cell uses the
strip bounds of \S\ref{sec:lap012} together with $|g(V\vp_\eta)-g(V)|\le2\varepsilon_g(V_T)$ and
$|g(V)|\ge1-\varepsilon_g(V_T)$, where $\varepsilon_g(11.99)=1.787\cdot10^{-5}$; since $\varepsilon_g$ decreases, these
bounds hold for every $V\ge11.99$. The tail cell overlaps the last tile on $[11.99,12]$.
\end{enumerate}
The certified relative error is therefore at most $0.19$.
\cert{lap34\_cert.py 2.0 12.0 0.02; lap34\_cert.py tail 11.99}{third/logs/lap34\_cert.log}
\end{lemma}

\emph{Validation (not part of the proof).} \texttt{lap34\_validate.py}: $70$ exact points at $n=150,300,821,2000$, with
$\mu$ up to $0.1050$ and $\lambda$ down to $4.0$; the true $|c/\Amp-g|$ lies in $1.8\cdot10^{-4}$--$1.6\cdot10^{-2}$, which
is $24$--$435\times$ below the certified bound, and every sign is negative. Two audits checked a further $36$ and $31$
points with independent exact code (the second by a multi-modular CRT method that shares no code with the others), up to
$n=5000$ and $V=46.8$, including points straddling the tile/tail junction: every sign negative, true error $28$--$579\times$
below the certified bound.

\subsubsection{Coverage lemma for the saddle}\label{sec:t-sadcov}
Write $A(V):=\alpha(V)/V^2$. Then $A$ is continuous, $A\to1/5$ as $V\to0$, and $A\to0$ as $V\to\infty$; and $\alpha$ is
increasing, since $\alpha'=Vq(V)/5$ with $q=1/(e^V-1)-5/(e^{5V}-1)>0$.

\begin{lemma}\label{lem:t-sadcov}
$A\ge0.04424$ on $[0.001,2]$ ($5858$ tiles, using $A\ge\alpha(v_1)/v_2^2$) and $A\ge0.1998$ on $(0,0.001]$.
\cert{cls012\_cov.py}{third/logs/cls012\_cov.log}
\end{lemma}

For classes $0$--$2$ the saddle equation is $A(V)=(m-1/6)/(25n^2)<0.4\mu$; for classes $3$--$4$ it is
$A(V)=(m-1/6-5n)/(25n^2)$, which is smaller still. When $\mu<0.105$, the right-hand side is $<0.042$. By the intermediate
value theorem a root exists, and every root has $V>2$. The representation is exact for every $W>0$, so any root may be
used.

%---------------------------------------------------------------------
\subsection{The centre, \texorpdfstring{$\mu\ge0.105$}{mu >= 0.105}}\label{sec:t-centre}
%---------------------------------------------------------------------

Throughout this subsection $n\ge821$ and $0.105\le\mu\le\tfrac12$. We use $r=e^{-L/(5n)}$ and
$q=\omega_5^he^{-\sigma/(5n)}$, with $\sigma=L-iy$, so that $y=5n(\theta-2\pi h/5)$;
$\Phi_c(s)=\log(1-\omega_5^{hc}e^{-s})$ (principal branch) and $Q=\sum_{c=1}^4\Phi_c=\log\frac{1-e^{-5s}}{1-e^{-s}}$;
$\Lambda(\sigma)=\int_0^1Q(t\sigma)\,dt$; and $C_h(\sigma)=\sum_c(c/5-\tfrac12)(\Phi_c(\sigma)-\Phi_c(0))$. The starting
point is Cauchy's formula for the polynomial $P_n$:
\[
c(m)=\frac1{2\pi}\int_{-\pi}^{\pi}P_n(re^{i\theta})\,r^{-m}e^{-im\theta}\,d\theta,\qquad r^{-m}=e^{2n\mu L}.
\]

\subsubsection{Near-peak piece (A): \texorpdfstring{$|y|\le1.1$}{|y| <= 1.1} around \texorpdfstring{$\theta=2\pi h/5$}{theta = 2 pi h/5}, \texorpdfstring{$h=1,\dots,4$}{h = 1..4}}\label{sec:nearpeak}

\begin{lemma}[finite Euler--Maclaurin]\label{lem:t-fem}
For $0\le a<1$,
\[
\sum_{t=0}^{n-1}f(t+a)=\int_0^nf+B_1(a)(f(n)-f(0))+\frac{B_2(a)}2(f'(n)-f'(0))-\int_0^n\frac{\tilde B_2(x-a)}2f''(x)\,dx .
\]
\end{lemma}
\begin{proof}
Apply the Euler--Maclaurin formula of order $2$ with offset $a$ on each $[t,t+1]$ and sum over $t$; the argument is the same
as for Lemma~\ref{lem:t-em}, but on a finite range.
\end{proof}
The constants are $B_2(1/5)=B_2(4/5)=1/150$ and $B_2(2/5)=B_2(3/5)=-11/150$. Since $\tilde B_2\in[-1/12,1/6]$, the last
term is at most $\frac1{12}\int|f''|$. The identity was checked numerically to $10^{-15}$.

\begin{lemma}\label{lem:t-logS}
Apply Lemma~\ref{lem:t-fem} to $f(x)=\Phi_c(x\sigma/n)$ with $a=c/5$, and sum over $c$. The terms $k=5t+c$, $t<n$, are
exactly the factors of $P_n$. This gives
\[
\big|\log P_n(q)-n\Lambda(\sigma)-C_h(\sigma)\big|\le\frac{K_h(\sigma)}n,
\]
where
\[
K_h=\sum_c\Big[\frac{|B_2(c/5)|}2|\sigma|\big(|\Phi_c'(\sigma)|+|\Phi_c'(0)|\big)+\frac{|\sigma|^2}{12}\sup_{t\in[0,1]}|\Phi_c''(t\sigma)|\Big],
\]
valid for $|y|<2\pi/5$ and $L\ge0$, uniformly down to $L=0$. The reason is that $\Phi_c$ has no singularity on the segment
$[0,\sigma]$, because $\omega_5^{hc}\ne1$.
\end{lemma}

Two exact facts are used: $|e^{C_h(L)}|=1$ for real $L$ (enclosed in the certificate, not assumed), and the four peaks have
equal modulus exactly. The class margin is then
\[
M_s(L)=\sgn_s\Re\sum_h\omega_5^{-hs}e^{C_h(L)},\qquad \sgn_0=+1,\ \sgn_{1,\dots,4}=-1 .
\]

\emph{Saddle.} $\mu=-\Lambda'(L)/2$. Since $Q'(0)=\sum_c\omega_5^c/(1-\omega_5^c)=-2$, we get $\Lambda'(0)=-1$ and
$\mu(0)=1/2$.

\begin{lemma}[saddle endpoint]\label{lem:t-saddle}
$\mu(2.1)<0.105$, by the real mean-value theorem on $4096$ $t$-pieces: the log prints
\texttt{mu(2.1) = -p'/2 in [0.104 +/- 3.02e-4] < 0.105: True}.
\cert{centre3\_saddle\_end.py}{third/logs/centre3\_saddle\_end.log}
\end{lemma}
(An audit computes $\mu(2.1)=0.103772$ from the closed form, so $L(0.105)=2.0814$.) The near-peak certificate asserts
$a_L:=\Lambda''(L)>0$ on every $L$-box, so $\mu(L)$ is strictly decreasing. Therefore every $\mu\in[0.105,1/2]$ has a
unique saddle $L\in[0,2.1]$.

\emph{Laplace step.} Put
\[
\cN(n,L)=\frac{e^{n\Lambda(L)+2n\mu L}\sqrt{2\pi/(na_L)}}{10\pi n}.
\]
Then
\[
c(m)=\cN(n,L)\Big[\sum_{h=1}^4\omega_5^{-hm}e^{C_h(L)}(1+\rho_h)+\mathrm{strip}+\mathrm{OFF}\Big],
\]
using $d\theta=dy/(5n)$; the linear phase cancels at the saddle, since $-in\Lambda'(L)y=2in\mu y$; and each peak contributes
$\cN\omega_5^{-hm}e^{C_h}$. The three error pieces are bounded as follows.
\begin{itemize}
\item \emph{Inner window $|y|\le0.35$.} Write $\psi(y)=\Lambda(L-iy)-\Lambda(L)+i\Lambda'(L)y$ and
$Z=n\psi+na_Ly^2/2+\Delta C$. The odd part $T=n\psi_3y^3+C_1y$ integrates to exactly $0$ against the Gaussian. The rest is
bounded using $|e^Z-1-Z|\le\frac{|Z|^2}2\max(1,e^{\Re Z})$; $\Re\psi\le-b y^2/2$, where $b$ is a certified lower bound for
$\Re \Lambda''$ on the window; Taylor constants from $\Lambda^{(k)}(\sigma)=\int_0^1t^kQ^{(k)}(t\sigma)\,dt$, enclosed on
$t$-pieces; the factor $e^{K/n}-1$ and the Gaussian tail.
\item \emph{Strip $0.35\le|y|\le1.1$.} Lemma~\ref{lem:t-logS} is combined with
$\Re\psi(y)=-\int_0^{|y|}(|y|-u)\Re \Lambda''(L-iu)\,du$, with $\Re \Lambda''$ bounded below on boxes. This gives terms of
the form $\sqrt n\,e^{-nD}$, which are non-increasing for $n\ge1/(2D)$; that condition is asserted. The strip and the
Gaussian tail are normalised with the upper curvature $a_{up}$, which is the safe direction.
\item \emph{Uniformity in $n$.} $n$ enters only through $K/n$, the $1/n$ moment terms, the Gaussian tail
$e^{-nAY^2/2}/\sqrt n$, and the strip. Each of these is non-increasing in $n$. So a bound at $n=N_0=821$ holds for all
$n\ge821$.
\end{itemize}

\begin{lemma}[near peak]\label{lem:t-nearpeak}
On $210$ $L$-boxes of width $0.01$ covering $[0,2.1]$, each with $28$ inner $y$-boxes, $60$ strip $y$-boxes and $64$
$t$-pieces, every quantity enclosed over its full box, at $N_0=821$: for every box and every class,
$M_s>\sum_h|e^{C_h}|(\rho_h+\mathrm{strip}_h)=:\mathrm{err}_s$. The worst ratios $\mathrm{err}/\mathrm{margin}$ are
$0.0566$, $0.2225$, $0.2241$, $0.7977$, $0.8469$ for the classes $0,1,2,3,4$. The minimum slack
$\min_s(M_{s,lo}-\mathrm{err}_{s,up})$ is $0.018703$, attained for class $4$ on the box $L\in[2.09,2.10]$.
\cert{centre3\_peak.py 821 210 64 (reproduced inside centre3\_comb.py)}{third/logs/centre3\_peak\_821.log, centre3\_comb\_821.log}
\end{lemma}
Since $L(0.105)=2.0814$, only the boxes up to $[2.08,2.09]$ are needed; the last box $[2.09,2.10]$ is certified as well.

\subsubsection{The off-peak lemma (B)}\label{sec:offpeak}

\begin{lemma}[off-peak lemma]\label{lem:offpeak}
Let $n\ge821$ and $L\in[0,2.1]$, and let $\theta$ satisfy $|5n(\theta-2\pi h/5)|>1.1$ for $h=1,\dots,4$, taken mod $2\pi$.
Then
\[
\log|P_n(e^{-L/(5n)}e^{i\theta})|-n\Lambda(L)\le-\mathcal E(n),\qquad \mathcal E(n)=1.5\log n+6.3\quad(\mathcal E(821)=16.36).
\]
\end{lemma}
Every certificate uses the constant $6.3$ as the exact ball \texttt{arb('6.3')}. The proof occupies the rest of this
subsection: the $\theta$-circle is covered by the pieces (E), (E0), (N) and (F) below, each certified.

\emph{Closure.} For fixed $n$, $\log|P_n(e^{-L/(5n)}e^{i\theta})|$ is continuous in $(\theta,L)$, with values in
$[-\infty,\infty)$. So each bound ``$\le n\Lambda(L)-\mathcal E(n)$'' certified on a box also holds on the closure of that
box. This closes the float-endpoint slivers $|y|\in(1.1,\mathrm{fl}(1.1))$ at the start of (E) and $(\mathrm{fl}(\pi),\pi]$
at the end of (F) (the Arb run of (F) also closes the latter by ending at \texttt{nextafter(pi, 4)}). The sliver
$L\in[0.1,\mathrm{fl}(0.1))$ at the lower end of the $L$-boxes of (E), (E0) and (N), where
$\mathrm{fl}(0.1)=0.1+5.5\cdot10^{-18}$, is closed directly: the lowest $L$-box of these pieces starts at $0.0999$, below the
exact $0.1$, so $L=0.1$ is inside a certified box and the reduction of Lemma~\ref{lem:monoL} from $L<0.1$ to $L'=0.1$ costs
exactly $n(0.1-L)\le 0.1\,n$, which is what each piece charges (\texttt{arb('0.1')}).

Since $P_n$ has real coefficients, $|P_n(\bar q)|=|P_n(q)|$, so it suffices to take $\theta\in[0,\pi]$. The pieces (E) and
(N) use both signs of $y$ near $h=1,2$, which also covers $h=3,4$.

\emph{Closed forms.} For $h\ne0$,
$\Lambda(\sigma)=\big(\Li_2(e^{-5\sigma})/5-\Li_2(e^{-\sigma})+4\zeta(2)/5\big)/\sigma$, using
$\sum_{c=0}^4\Li_2(\omega_5^cx)=\Li_2(x^5)/5$. For $h=0$, $\Lambda_0(\sigma)=4(\Li_2(e^{-\sigma})-\zeta(2))/\sigma$. $Q$ is
decreasing and convex with $|Q'|\le2$: indeed $Q''=[g_2(s/2)-g_2(5s/2)]/s^2$ with $g_2(x)=x^2/\sinh^2x$ decreasing. Since
$\Lambda'=(Q-\Lambda)/L$, $\Lambda$ is decreasing on $L>0$.

\begin{lemma}[monotonicity in $L$]\label{lem:monoL}
For $0\le L\le L'$,
$\log|P_n(re^{i\theta})|\le\log|P_n(r'e^{i\theta})|+n(L'-L)$, where $r=e^{-L/(5n)}$ and $r'=e^{-L'/(5n)}$.
\end{lemma}
\begin{proof}
We have $|1-\varrho e^{ia}|^2=(1-\varrho)^2+4\varrho\sin^2(a/2)$. For $\varrho'\le\varrho\le1$, this gives
$|1-\varrho e^{ia}|^2\le(\varrho/\varrho')|1-\varrho'e^{ia}|^2$. Take $\varrho=r^k$, $\varrho'=r'^k$, and use
$\sum_{k\le5n,\,5\nmid k}k=10n^2$: the total factor on $\log|P_n|$ is $\frac12\cdot\frac{L'-L}{5n}\cdot10n^2=n(L'-L)$.
\end{proof}

Pieces (E), (E0) and (N) are certified on $L\in[0.1,2.1]$. For $L<0.1$ they use Lemma~\ref{lem:monoL} at $L'=0.1$ together
with $n\Lambda(L)\ge n\Lambda(0.1)$ ($\Lambda$ decreasing, and $y$ unchanged because $\theta$ is fixed). The cost is $0.1n$,
which the boxes touching $L=0.1$ are required to absorb.

\paragraph{Piece (E): $h\in\{1,2\}$, $1.1\le|y|\le6$.} $\Phi_c$ is analytic on $\Re s>0$ and at $s=0$. So
Lemma~\ref{lem:t-fem} applies for every $y$, even past $|\sigma|=2\pi/5$; the disc was needed only for the Taylor step in
\S\ref{sec:nearpeak}. This gives $\log|P_n|\le n\Re \Lambda(\sigma)+\Re C_h(\sigma)+K_h(\sigma)/n$, with the remainder in
integral form. The certificate uses $20000$ adaptive boxes. On each box it checks
$\mathcal V(N_0)=-N_0g_E+\Re C+K/N_0+\mathcal E(N_0)<0$, where $g_E=\Lambda(L_{hi})-\Re \Lambda(\sigma)-0.1\chi$, and
$g_E>1.5/N_0$, so that $\mathcal V(n)$ is non-increasing in $n$. The largest $K$ is $5.6\cdot10^3$, near a singular
crossing at $L=0.1$. The logged worst values of $\mathcal V(N_0)$ are $-2.49$, $-2.27$ ($h=1$, $y\gtrless0$) and
$-2.55$, $-2.18$ ($h=2$).
% Earlier drafts quoted -12.03; that predates extending the lowest L-box to 0.0999 (which widens it and loosens the bound).
\cert{offpeak\_em.py 821 6.0 32}{third/logs/offpeak\_em\_821.log}

\paragraph{Piece (E0): $h=0$, $|y|\le20$.} Here $\Phi=\log(1-e^{-s})$ is singular at $0$, so the $t=0$ terms are split
off and Lemma~\ref{lem:t-fem} is applied on $t=1,\dots,n-1$; note $\sum_cB_1(c/5)=0$. This gives
\[
\log|P_n|\le n\Re \Lambda_0(\sigma)+\Psi_0(\varsigma)+0.08\big(|\varsigma\Phi'(\varsigma)|+|\varsigma||\Phi'(\sigma)|\big)
+\frac{|\varsigma||\sigma|}3\int_{1/n}^1|\Phi''(t\sigma)|\,dt,
\]
where $\varsigma=\sigma/n$,
$\Psi_0(\varsigma)=4+\log(24/625)+\sum_c\Re\ell(c\varsigma/5)-4\Re[\varsigma^{-1}\int_0^\varsigma\ell]$ and
$\ell(s)=\log((1-e^{-s})/s)$. The $\log|\varsigma|$ terms cancel exactly. For $\Re s\ge0$ and $|s|\le0.1$:
$|(1-e^{-s})/s-1|\le|s|e^{|s|}/2$, hence $|\ell(s)|\le|s|$; $|s\Phi'(s)|\le1/(1-|s|e^{|s|}/2)$; and
$|\Phi''(s)|\le1/(|s|^2m_0^2)$, with $m_0=1-|s|e^{|s|}/2$. Consequently $\Psi_0\le0.7403+6|\sigma|/N_0$. The certificate uses
$4945$ boxes, and the worst value is $\mathcal V=-0.23$; this is a bisection artefact, as the true slack is in the hundreds.
\cert{offpeak\_em0.py 821 20 32}{third/logs/offpeak\_em0\_821.log}

\paragraph{Piece (N): $6\le|y|\le0.25n$ ($20\le y$ for $h=0$).} Since $|q|<1$,
\[
\log|P_n|=-\Re\sum_jG_j/j\le\sum_j|G_j|/j,\qquad G_j=\frac{(1-e^{-j\sigma})A_j}{1-e^{-j\varsigma}},\qquad
A_j=\sum_c\omega_5^{hcj}e^{-cj\varsigma/5},
\]
with $\varsigma=\sigma/n$; $A_j$ is the sum over the four $c$ for fixed $j$ (the factor $1-e^{-j\sigma}$ comes from summing
the geometric series over $t<n$). The bounds on $A_j$ are $|A_j|\le\min(4,\bar A_j+2j|\varsigma|)$ and
$|A_j|\le\alpha_N(jL/n)$ with $\alpha_N(a)=\sum_ce^{-ca/5}$, where $\bar A_j=|\sum_c\omega_5^{hcj}|$ is $1$ or $4$.
\begin{itemize}
\item \emph{Part I, $j\le\pi n/y$.} Use $1/|1-e^{-z}|\le1/|z|+c_1$ with $c_1=1+1/\pi$. This is the maximum modulus of
$1/(1-e^{-z})-1/z$ on the rectangle $[0,20]\times[-\pi,\pi]$: it is at most $0.593$ on the left edge, $1+1/\pi$ on the top
and bottom edges, and $1.06$ on the right edge. Every $z=j\varsigma$ in Part I lies in the rectangle, since
$|\Im z|=jy/n\le\pi$ and $\Re z=jL/n\le\pi L/y\le1.1$. This gives Part I $\le nM_N(\sigma)+R_1$, where
$M_N(\sigma)=\sum_j|1-e^{-j\sigma}|\bar A_j/(j^2|\sigma|)$ and $R_1\le B_N\big(1+\log(\pi n/y)\big)+O(1)$, with
$B_N=2+1.6c_1$ for $h\ne0$ and $B_N=2+4c_1$ for $h=0$.
\item \emph{Part II, $j>\pi n/y$.} Split the $j$ into periods $k\ge1$, with $jy/n\in((2k-1)\pi,(2k+1)\pi]$. Write
$a=jL/n$ and $b=jy/n$, so that $j\varsigma=a-ib$. On period $k$: $1/j\le y/((2k-1)\pi n)$;
$|1-e^{-j\varsigma}|\ge\max\big(1-e^{-a},\ 2e^{-a/2}|\sin(b/2)|\big)$; $|\sin(b/2)|\ge|b-2\pi k|/\pi$ on the period. The two
$j$ nearest the pole $b=2\pi k$ are bounded by $G_k=1/(1-e^{-a_k^-})$, where $a_k^-$ is the smallest $a$ on the period. The
other $j$ on the period are bounded through $|\sin(b/2)|$ and summed with
$\sum_{i\le M'}\min(G,D/i)\le\min\big(M'G,\ D(1+\log^+(2M'G/D))\big)$. Consequently: summed over $k$, the pole terms total
at most $\beta u^2n/L$, with $u=y/n$ and $\beta=1+e^{-\pi L/u}$; the remaining terms of the periods total $O(\log^2n)$; the
periods with $a_k^->2$ give a constant tail, bounded with $\int_2^\infty\phi(a)/a\,da$ in closed form. Part II is
non-increasing in $n$ at fixed $(L,y)$: $M'/n$, $D/n$ and $\log^+(2M'G/D)$ do not increase, and $\beta_k$ decreases. (N)
stops at $u=0.25$ because the pole term $\beta u^2/L$ must stay below $\Lambda(L)-0.1$ at $L=0.1$. The triangle-bound
majorant blows up at low-denominator rationals such as $\theta=\pi$ and $2\pi/3$, and that is where (F) takes over.
\end{itemize}
The certificate has two parts: \textbf{(N-a)} $y\in[Y_E,200]$ on $(L,y)$-boxes, at $N_0$, with the exact $M_N(\sigma)$
($400$ terms plus a tail) and slope check $g\ge(B_N+1.5)/N_0$ (monotonicity in $n$ follows from $R_1=A+B_N\log n$, Part II
being non-increasing, and the slope condition); \textbf{(N-b)} $y\in[200,0.25n]$ on $(L,u)$-boxes, in the form
$-nG+a+b\log n+c\log^2n$ with $G\ge(b+2c\log N_0)/N_0$. The worst values are $-147.9/-23.3$ for $h=1$ and $-71.3/-1.8$
for $h=0$. The script runs only $h\in\{0,1\}$ with $y>0$. This suffices: the majorant depends on $h$ only through
$\bar A_j$, which equals $1$ if $5\nmid j$ and $4$ if $5\mid j$ for every $h=1,\dots,4$, and it depends on $y$ only through
$|1-e^{-j\sigma}|$, $|\sigma|$ and $|y|$, all invariant under $y\mapsto-y$.
\cert{offpeak\_tb.py 821 0.25}{third/logs/offpeak\_tb\_821.log}

\paragraph{Piece (F): $\operatorname{dist}(\theta,\{0,2\pi/5,4\pi/5\})\ge0.0499$, $\theta\le\pi$, $L\in[0,2.1]$.}
The chunk length is $T=40$ on the bands $0.05\le\operatorname{dist}<0.1$ and $T=20$ elsewhere. Put $M_T=\lfloor n/T\rfloor$,
$\lambda_T=TL/n$, $\mu_r=L/(5n)$ and $\Pi(Z)=\prod_{k'\le5T,\,5\nmid k'}(1-Ze^{ik'\tilde\theta})$ with
$\tilde\theta=\theta+i\mu_r$. Here the radial drift $r^{k'}$ appears as an imaginary shift of $\theta$.

\begin{lemma}[chunk lemma]\label{lem:chunk}
\[
\log|P_n|\le\sum_{m'<M_T}\cB(m'\lambda_T)+4(T-1)\log(1+e^{-(L-\lambda_T)}),\qquad\cB(s)=\max_{|Z|=e^{-s}}\log|\Pi(Z)|,
\]
\[
n\Lambda(L)\ge\sum_{m'<M_T}TQ(m'\lambda_T)-T\big[\lambda_T+(\log5-Q(L))/2\big].
\]
\end{lemma}
\begin{proof}
Write $n=M_TT+r_T$ with $0\le r_T<T$. For $m'<M_T$, the factors with $5Tm'<k\le5T(m'+1)$ form one chunk: with $k=5Tm'+k'$
and $Z_{m'}=q^{5Tm'}$, the chunk equals $\Pi(Z_{m'})$ with $|Z_{m'}|=e^{-m'\lambda_T}$, so its log-modulus is at most
$\cB(m'\lambda_T)$. This holds for every $n$, whether or not $T\mid n$. The factors with $5TM_T<k\le5n$ and $5\nmid k$ are
left over; there are $4r_T\le4(T-1)$ of them. Each satisfies $|1-q^k|\le1+e^{-kL/(5n)}$, and $kL/(5n)>TM_TL/n>L-\lambda_T$
because $TM_T>n-T$. So the leftover factors cost at most $4(T-1)\log(1+e^{-(L-\lambda_T)})$. The second line uses the
tangent line of $Q$ at each $m'\lambda_T$, convexity telescoping $-Q'(m'\lambda_T)\lambda_T\le Q((m'-1)\lambda_T)-Q(m'\lambda_T)$,
$|Q'|\le2$ for the $m'=0$ term, $Q>0$, and $M_T\lambda_T\le L$.
\end{proof}

\emph{Reductions.} $\cB$ is convex and non-increasing in $s$ (Hadamard three circles). So on an $s$-grid of step $0.1$,
$\cB$ is below its chord and $TQ$ is above its tangent, which gives cell lower bounds $d_i$. $\sup_Z\log|\Pi|$ is
subharmonic in $\tilde\theta$, since $\Pi$ is entire in $\tilde\theta$; so on each rectangle
$[\theta_a,\theta_b]\times[0,\mu_{\max}]$, with $\mu_{\max}=2.1/(5N_0)$, the maximum is attained on the boundary, and the
drift needs no interior treatment. The assembly uses left Riemann sums of the running minimum of $d_i$; for small $L$ it
uses $M_Td_{\min}$ instead. It is checked on the $L$-intervals $[i/100,(i+1)/100]$, $i=0,\dots,209$, which cover $[0,2.1]$
including $L=0$, at $N_0$, with the slope condition, so it holds for all $n\ge N_0$.

\begin{lemma}[piece (F)]\label{lem:pieceF}
On $163$ rectangles and $22$ $s$-values, in Arb ball arithmetic over dyadic cells, the assembled margin is positive on every
rectangle; the minimum is $47.91$ (piece $7$, rectangle $26$, ending at $\pi^+$, $T=20$). The tightest $T=40$ band has margin
$51.8$.
\cert{offpeak\_chunk\_arb.py 821 \$p (p=0,\dots,7), then offpeak\_chunk\_arb.py summarize}{third/logs/arb\_F\_summary\_full.log (merging the 13 logs arb\_F\_p*.log)}
\end{lemma}
The same bound was first certified in float64 interval arithmetic with directed-rounding margins (only correctly rounded
$+,-,\times,\div$ and exact \texttt{frexp}/\texttt{ldexp}; transcendental inputs from Arb with explicit margins),
with identical minimum margin $47.9$ (\texttt{offpeak\_chunk.py 821 all}, log \texttt{offpeak\_chunk\_821.log}); the Arb
certificate supersedes it and also fixes the endpoint $\pi$ and an inexact key in the assembly. An early exit in the Arb
cell test that could only reject cells, never accept them, stalled the first runs; it was removed, and every rectangle
logged before the fix remains valid (an audit re-ran $11$ of them, including the worst, with identical margins).

\emph{Coverage of the $\theta$-circle.} The pieces meet or overlap at every handover: (A)--(E) abut at $|y|=1.1$ (the
statement has a strict $>$); (E)--(N) abut at $|y|=6$; (E0)--(N) abut at $y=20$; (N-a)--(N-b) meet at $y=200$, which needs
$0.25n\ge200$, i.e.\ $n\ge800$; (N)--(F): (N) reaches $|\theta-2\pi h/5|=0.25n/(5n)=0.05$ for every $n$, and (F) starts at
$0.0499$, an overlap of $10^{-4}$ rad. The root $6\pi/5$ is $0.63$ away from $\pi$, so $[0,\pi]$ needs no exclusion around it.
The (F) rectangles tile each piece exactly, and the pieces abut exactly. The float-endpoint slivers are closed by the
closure remark. This proves Lemma~\ref{lem:offpeak}.

\emph{Spot validation (not part of the proof).} $312$ random points and, in a second audit, $5570$ adversarial points (for
$n$ from $821$ to $10^5$, window edges, box endpoints, $L$ down to $10^{-9}$, located local maxima) were evaluated with
exact Arb full products: $0$ violations; the smallest true slack is $43.18$, at $n=821$, $L=2.1$, $|y|=1.1^+$.

\subsubsection{Combination}\label{sec:comb}
On the off-peak set, Lemma~\ref{lem:offpeak} gives $|P_nr^{-m}|\le e^{n\Lambda(L)+2n\mu L-\mathcal E(n)}$. The set has
measure at most $2\pi$, so
\[
|\mathrm{OFF}|\le\frac{1}{\cN}\cdot\frac{2\pi}{2\pi}e^{n\Lambda+2n\mu L-\mathcal E(n)}=10\pi n\sqrt{\frac{na_L}{2\pi}}\,e^{-\mathcal E(n)}
=10\pi\sqrt{\frac{a_L}{2\pi}}\,e^{-6.3}.
\]
The powers of $n$ cancel exactly, so this bound is independent of $n$. It grows with $a_L=\Lambda''(L)$. The largest
$a_L$ occurs at $L=0$, where $a_L=2/3$, while the minimum slack of Lemma~\ref{lem:t-nearpeak} occurs at $L\approx2.1$; so
the comparison must be made per $L$-box.

\begin{lemma}[combination]\label{lem:comb}
On every one of the $210$ $L$-boxes and for every class $s$, with
$\mathrm{OFF}_{up}=10\pi\sqrt{a_{up}/2\pi}\,e^{-6.3}$ computed in Arb with the exact constant \texttt{arb('-6.3')}, the
program asserts $M_{s,lo}-\mathrm{err}_{s,up}-\mathrm{OFF}_{up}>0$ in Arb. Over all boxes and classes the minimum of
slack minus OFF is $0.010043$, on the box $[2.09,2.10]$ (slack $0.018703$, OFF $0.008660$); on the last box actually
needed, $[2.08,2.09]$, it is $0.0120$ ($0.020699-0.008696$).
\cert{centre3\_comb.py 821 210 64 0 105; centre3\_comb.py 821 210 64 105 210}{third/logs/centre3\_comb\_821.log}
\end{lemma}

\begin{corollary}\label{cor:centre}
For $n\ge821$, $\mu\in[0.105,\tfrac12]$ and every class, $c(m)$ has the conjectured sign.
\end{corollary}
\begin{proof}
Lemmas~\ref{lem:t-saddle}, \ref{lem:t-nearpeak}, \ref{lem:offpeak} and~\ref{lem:comb}.
\end{proof}

\emph{Validation (not part of the proof).} An audit evaluated $c(m)/\cN$ by quadrature on the windows, validated against
exact coefficients at $n=300$ (35 points, agreement to $4\cdot10^{-11}$), and then checked $120$ points at $n=821$, $857$
and $900$, with $\mu\in\{0.105,0.15,0.2,0.3,0.4,0.45,0.49,0.5\}$ and all five classes: every sign correct; realised error at
most $1.65\cdot10^{-3}$ (at most $4.7\cdot10^{-4}$ in classes $3$ and $4$), against a certified $\mathrm{err}$ of
$0.06$--$0.21$.

%---------------------------------------------------------------------
\subsection{Coverage and proof of Theorem~\ref{thm:third-main}}\label{sec:t-coverage}
%---------------------------------------------------------------------

\begin{theorem}[coverage]\label{thm:t-coverage}
Let $n\ge980$ and $0\le m\le5n^2$. Then $(n,m)$ is covered by at least one of the following: the band certificates of
\S\ref{sec:t-band} (if $\lambda<4$); the Laplace certificates of \S\ref{sec:t-laplace} (if $\lambda\ge4$ and $\mu<0.105$);
the centre certificate of \S\ref{sec:t-centre} (if $\mu\ge0.105$). Each certificate that applies proves the conjectured
sign of $c_n(m)$.
\end{theorem}
\begin{proof}
The three conditions partition $\{0\le m\le5n^2\}$. ($\lambda\ge4$ is automatic when $\mu\ge0.105$, since
$m\ge1.05n^2>4(5n+1)$ for $n\ge22$, but it is not needed.) It remains to check that each certificate's hypotheses hold on
its whole region.
\begin{enumerate}
\item \emph{Band, $m<4N$.} Classes $3,4$: Corollary~\ref{cor:t-block} covers $m<2N$ for every $n$; for $2N\le m<4N$, the
exact computation covers $n\le2000$ and $\Psi$ covers $n>2000$. Classes $0,1,2$: the exact computation covers $n\le2000$,
and Lemma~\ref{lem:t-band012} covers $n>2000$. The only threshold is $n\ge291$ for the exact band run, which holds.
\item \emph{Laplace, $\lambda\ge4$, $\mu<0.105$.} By \S\ref{sec:t-sadcov} the saddle equation, for either class family,
has a root, and every root satisfies $V>2$. Fix a root. Then $W=V/(5n)$, the representation of \S\ref{sec:t-repr} is
exact, and $X=5n/V$. If $V\in[2,12]$, $V$ lies in one of the $500$ tiles $[V_{lo},V_{hi}]$; for both class families the
tiles are contiguous and the last tile is snapped to end exactly at $12$. Classes $0$--$2$: $X\ge1455/V_{hi}$, since
$n\ge291$, and $X\ge4V_{lo}/a_5$, since $\lambda\ge4$; so $X\ge X_{\min}$ of the tile. Classes $3$--$4$:
$X\ge4105/V_{hi}$, since $n\ge821$, and $X>3V_{lo}/a_5$, since $\lambda\ge4$. If $V\ge11.99$, the tail cell applies; for
both class families it is run from $V_T=11.99$, so it overlaps the last tile. It needs $X\ge4V/a_5$ (classes $0$--$2$) or
$X\ge X^*=216.3$ together with $V\le a_5X/3$ (classes $3$--$4$), and these follow from the same two inequalities. Every
certificate term is non-increasing in $X$, so the tile's certified inequality holds at this $X$. The thresholds are
$n\ge291$ (classes $0$--$2$) and $n\ge821$ (classes $3$--$4$), both below $980$.
\item \emph{Centre, $\mu\ge0.105$.} Since $m\le5n^2$, $\mu\le\tfrac12$. By \S\ref{sec:nearpeak}, $\mu\mapsto L$ is a
bijection from $[0.105,\tfrac12]$ onto $[0,L(0.105)]\subset[0,2.09]$ (since $L(0.105)=2.0814$), and every $L$-box in
$[0,2.1]$ is certified at $N_0=821$ for all $n\ge821$. The integration circle is split into the four near-peak windows
$|y_h|\le1.1$ and their complement. The complement is covered by (E), (E0), (N) and (F), with the handovers listed in
\S\ref{sec:offpeak} and the closure remark for the float endpoints, and Lemma~\ref{lem:offpeak} holds there. The per-box
combination of Lemma~\ref{lem:comb} then gives the sign. The threshold is $n\ge821$.
\end{enumerate}
So all thresholds ($291$, $821$, and $2000$ for the switch between exact and analytic in the band) are satisfied or bridged
for $n\ge980$. The boundaries $\lambda=4$ and $\mu=0.105$ are assigned to the Laplace region and the centre respectively.
Each certificate includes its closed boundary: the Laplace certificates use $\lambda\ge4$, and the centre uses
$\mu\ge0.105$ through $\mu(2.1)<0.105$. So nothing is left out.
\end{proof}

\begin{proof}[Proof of Theorem~\ref{thm:third-main}]
By Lemma~\ref{lem:palin} it suffices to take $0\le m\le5n^2$. For $n\le979$ the result is Theorem~\ref{thm:t-exact}. For
$n\ge980$ it follows from Theorem~\ref{thm:t-coverage}, together with Corollaries~\ref{cor:t-band34}
and~\ref{cor:t-band012}, Lemmas~\ref{lem:lap012-tiles}, \ref{lem:lap012-tail} and~\ref{lem:lap34}, and
Corollary~\ref{cor:centre}. The analytic inputs are Theorem~\ref{thm:G5}, the elementary lemmas proved above, and the
standard facts of Appendix~\ref{app:facts}.
\end{proof}

%=====================================================================
\section{Computational verification}\label{sec:computations}
%=====================================================================

All certificates are Python programs using FLINT \cite{FLINT} and Arb \cite{Arb} through \texttt{python-flint}. They are
published with the paper \cite{Code} in two directories, \texttt{cubic/} (Theorem~A) and \texttt{third/} (Theorem~B), each
with its programs in \texttt{certificates/}, its logs in \texttt{logs/}, and a script \texttt{run\_all.sh} that regenerates
every log. A certificate consists of a program, the log it writes, and a line that the log must contain; every program
aborts on a failed assertion. A printed bound is an upper (or lower) end of a certified enclosure.

\subsection{Certificates for Theorem~A}
Programs in \texttt{cubic/certificates/}, logs in \texttt{cubic/logs/}. The whole set is regenerated by
\texttt{cubic/run\_all.sh} in about three hours on one core; the uniform Laplace grid parallelises over its $36$ chunks
(\texttt{JOBS=16 ./run\_all.sh}).

\begin{center}
\small
\renewcommand{\arraystretch}{1.12}
\begin{longtable}{@{}>{\raggedright\arraybackslash}p{0.25\textwidth}>{\raggedright\arraybackslash}p{0.24\textwidth}>{\raggedright\arraybackslash}p{0.44\textwidth}@{}}
\toprule
certifies & program(s) & log: required line\\
\midrule
\endhead
finite range $n\le243$ (Cert.~\ref{cert:c-finite}) & \texttt{exact\_small\_n.py 1 200}; \texttt{201 225}; \texttt{226 243} &
\texttt{01\_exact\_small\_n\_\{a,b,c\}.log}: \texttt{EXACT n in [..]: violations 0} in all three\\
band identity and $g(i)<0$, $i\le36000$ (Cert.~\ref{cert:c-band}) & \texttt{band\_lemma.py 36000} &
\texttt{02\_band\_lemma.log}: \texttt{(a) ... True}, \texttt{(b) ... True}\\
band tail $i\ge36000$ (Prop.~\ref{prop:c-band}) & \texttt{band\_tail\_rigorous.py} &
\texttt{02b\_band\_tail.log}: \texttt{BAND TAIL: g(i) < 0 for all i >= 36000}\\
constants of Thm~\ref{thm:expansion3} & \texttt{sz\_expansion\_audit.py} & \texttt{03\_sz\_audit.log}: \texttt{E\_con=50.150108}\\
Poisson bound (Cert.~\ref{cert:c-poisson}) & \texttt{poisson\_check.py} & \texttt{04a\_poisson.log}: \texttt{max |eps| ... <= 1.420}\\
uniform Laplace grid (Cert.~\ref{cert:c-laplace}) & \texttt{run\_uniform.py} ($36$ chunks) &
\texttt{uni\_*.log}: every line \texttt{UNRESOLVED 0}; worst $B_{\rm tot}=0.1169$\\
closing lemma (Prop.~\ref{prop:c-closing}) & \texttt{ray\_asymptotic2.py} &
\texttt{07\_closing\_lemma.log}: \texttt{CLOSING LEMMA (v >= 3000, all N): CERTIFIED}\\
centre tables (Cert.~\ref{cert:c-centre}) & \texttt{centre\_mid\_\allowbreak uniform\{,2,3\}.py}, \texttt{centre\_g3\_\allowbreak uniform\{,2,3\}.py},
\texttt{centre\_clow\_uniform.py} & \texttt{*\_table.txt} ($115$, $115$, $74$ lines): no assertion errors\\
coverage and centre combination (Prop.~\ref{prop:c-coverage}) & \texttt{coverage.py} &
\texttt{08\_coverage.log}: \texttt{PASS A}--\texttt{PASS G}, \texttt{COVERAGE COMPLETE}\\
\midrule
\multicolumn{3}{@{}l}{\emph{Consistency checks, not part of the proof}}\\
expansion bound, $i\le108000$ & \texttt{expansion\_check.py 36000} & \texttt{03b\_expansion\_check.log}: \texttt{max |rest|/bound = 0.8660}\\
Appendix~\ref{app:elementary} checks & \texttt{hand\_checks.py}, \texttt{koksma\_cauchy\_check.py} & \texttt{04b}, \texttt{04c}: ratios $\le1$\\
Prop.~\ref{prop:c-reduce} end to end & \texttt{identity\_check.py} & \texttt{04d\_identity.log}: rel.\ err.\ $<10^{-8}$ on every line\\
\bottomrule
\end{longtable}
\end{center}
\texttt{expansion\_check.py} uses \texttt{mpmath} floating point with adaptive precision; it corroborates
\S\ref{sec:expansion3} and is not part of the proof.

\subsection{Certificates for Theorem~B}
Programs in \texttt{third/certificates/} (together with the shared modules \texttt{offpeak\_common.py},
\texttt{cls012\_lib.py}, \texttt{lap34\_lib.py} and the exact coefficient table \texttt{g5\_coeffs.pkl}), logs in
\texttt{third/logs/}. Runtimes are those printed in the logs or recorded by the audits, on a 2-core machine.

\begin{center}
\small
\renewcommand{\arraystretch}{1.12}
\begin{longtable}{@{}>{\raggedright\arraybackslash}p{0.2\textwidth}>{\raggedright\arraybackslash}p{0.27\textwidth}>{\raggedright\arraybackslash}p{0.38\textwidth}>{\raggedright\arraybackslash}p{0.07\textwidth}@{}}
\toprule
certifies & command & log: required line & time\\
\midrule
\endhead
exact, $1\le n\le150$ (Thm~\ref{thm:t-exact}) & \texttt{centre3\_exact.py 1 150} & \texttt{exact\_1\_150.log}: \texttt{DONE n=1..150 total violations 0} & 4\,s\\
exact, $150\le n\le291$ & \texttt{centre3\_exact.py 150 291} & \texttt{exact\_150\_291.log}: \texttt{DONE n=150..291 total violations 0} & 34\,s\\
exact, $291\le n\le410$ & \texttt{centre3\_exact.py 291 410} & \texttt{centre3\_exact\_291\_410.log}: \texttt{DONE ... violations 0} & 114\,s\\
exact, $410\le n\le580$ & \texttt{centre3\_exact.py 410 580} & \texttt{centre3\_exact\_410\_580.log}: \texttt{DONE ... violations 0} & 419\,s\\
exact, $580\le n\le820$ & \texttt{centre3\_exact.py 580 820} & \texttt{centre3\_exact\_580\_820.log}: \texttt{DONE ... violations 0} & 1357\,s\\
exact, $820\le n\le920$ & \texttt{centre3\_exact.py 820 920} & \texttt{centre3\_exact\_820\_920.log}: \texttt{DONE ... violations 0} & 2825\,s\\
exact, $920\le n\le979$ & \texttt{centre3\_exact.py 920 1000} & \texttt{centre3\_exact\_920\_1000.log}: $60$ lines \texttt{n=920} \dots\ \texttt{n=979 wrong\_sign=0}, then \texttt{Killed} (about $6$\,GB needed) & 3786\,s\\
band, exact, $291\le n\le2000$, $m<4N$ & \texttt{bandfix\_exact.py 291 2000} & \texttt{bandfix\_exact.log}: \texttt{sign failures (any class, m < 4N): 0} & 192\,s\\
coefficient table; Lemma~\ref{lem:t-block-exact} & \texttt{bandfix\_coeffs.py} & \texttt{bandfix\_coeffs.log}: \texttt{g5\_coeffs.pkl agrees ...} and \texttt{ONE-PART BLOCK (exact): ...} & $<1$\,min\\
Lemma~\ref{lem:t-L1}, $\Psi$ & \texttt{bandfix\_analytic.py 10001} & \texttt{bandfix\_analytic.log}: \texttt{L1: ... <= [0.017972 +/- 1.12e-7]}; \texttt{BAND (analytic): ... < 0.1977 for all N >= 10001} & $<1$\,min\\
band classes $0$--$2$ (Lemma~\ref{lem:t-band012}) & \texttt{cls012\_band.py} & \texttt{cls012\_band.log}: \texttt{PART A: ... sign failures 0}; \texttt{PART B: ... Q <= [5.1390e-20 ...]} & seconds\\
Laplace $0$--$2$, tiles (Lemma~\ref{lem:lap012-tiles}) & \texttt{cls012\_lap.py 2.0 12.0 0.02} & \texttt{cls012\_lap\_rerun.log} l.~1: \texttt{500 tiles, 0 failures; ... 0: 3.0700, 1: 1.6879, 2: 0.8339} & 8\,s\\
Laplace $0$--$2$, tail (Lemma~\ref{lem:lap012-tail}) & \texttt{cls012\_lap.py tail 11.99} & \texttt{cls012\_lap\_rerun.log} l.~2: \texttt{tail cell V >= 11.99 : (\{0: (3.0950...} & seconds\\
saddle coverage (Lemma~\ref{lem:t-sadcov}) & \texttt{cls012\_cov.py} & \texttt{cls012\_cov.log}: \texttt{5858 tiles, min A = [0.04424...]} & seconds\\
Laplace $3$--$4$ (Lemma~\ref{lem:lap34}) & \texttt{lap34\_cert.py 2.0 12.0 0.02}; \texttt{lap34\_cert.py tail 11.99} & \texttt{lap34\_cert.log}: \texttt{500 tiles, 0 failures}; margins \texttt{0.8363}, \texttt{0.8361}; \texttt{tail cell V >= 11.99 : (\{3: (0.8145634857582285, ...} & 76\,s\\
saddle endpoint (Lemma~\ref{lem:t-saddle}) & \texttt{centre3\_saddle\_end.py} & \texttt{centre3\_saddle\_end.log}: \texttt{... < 0.105: True} & seconds\\
near peak and combination (Lemmas~\ref{lem:t-nearpeak}, \ref{lem:comb}) & \texttt{centre3\_comb.py 821 210 64 0 105} and \texttt{... 105 210} & \texttt{centre3\_comb\_821.log}: \texttt{ALL\_OK} on every box; twice \texttt{COMBINATION ASSERTED (slack > OFF, arb) on every box; ... True} & $\approx$11\,min\\
off-peak (E) & \texttt{offpeak\_em.py 821 6.0 32} & \texttt{offpeak\_em\_821.log}: \texttt{piece (E) CERTIFIED for all n >= N0} & 57\,s\\
off-peak (E0) & \texttt{offpeak\_em0.py 821 20 32} & \texttt{offpeak\_em0\_821.log}: \texttt{piece (E0) CERTIFIED for all n >= N0} & 6\,s\\
off-peak (N) & \texttt{offpeak\_tb.py 821 0.25} & \texttt{offpeak\_tb\_821.log}: \texttt{piece (N) CERTIFIED for all n >= N0} & 2\,s\\
off-peak (F), Arb (Lemma~\ref{lem:pieceF}) & \texttt{offpeak\_chunk\_arb.py 821 \$p -j 2}, $p=0,\dots,7$; then \texttt{... summarize} & \texttt{arb\_F\_summary\_full.log}: \texttt{(F) arb: CERTIFIED for all n >= N0 (all rects OK)} & $\approx$1.5\,h\\
off-peak (F), float (superseded) & \texttt{offpeak\_chunk.py 821 all} & \texttt{offpeak\_chunk\_821.log}: \texttt{(F) all: CERTIFIED} & 470\,s\\
\bottomrule
\end{longtable}
\end{center}
\texttt{centre3\_peak\_821.log} (\texttt{centre3\_peak.py 821 210 64}) contains the near-peak part alone and is
reproduced inside \texttt{centre3\_comb\_821.log}. The coefficient table \texttt{g5\_coeffs.pkl} is produced by
\texttt{g5\_coef.py}. The file \texttt{arb\_F\_summary.log} of the working directory, an incomplete earlier summary, is
not a certificate and is not shipped.

\subsection{Exact ranges}
\begin{itemize}
\item Theorem~A: $P_{3,n}^3$ exactly for $n\le243$; exact $s(i)$ and $g(i)<0$ for $i\le36000$; the band identity for
$n\le60$.
\item Theorem~B: $P_{5,n}$ exactly for $n\le979$ (all classes, all $m$); the band $m<4N$ exactly for $291\le n\le2000$;
exact $s(i)$ for $i\le200000$; the one-part block for $t\le39999$; Lemma~\ref{lem:t-band012}(A) for $m\le200000$.
\end{itemize}

\subsection{Reproducibility}
The software versions are Python~3.12, \texttt{python-flint 0.9.0}, \texttt{mpmath 1.3.0}, \texttt{numpy 2.4.4},
\texttt{scipy 1.17.1} and \texttt{sympy 1.14.0}. The repository's top-level \texttt{Makefile} has targets
\texttt{verify-cubic} and \texttt{verify-third}, which call the two \texttt{run\_all.sh} scripts. The expensive steps are the
cubic uniform Laplace grid and exact range (about three hours on one core in total), and, for Theorem~B, the exact range
$n\le979$ (about $2.4$ hours of CPU time, and about $6$\,GB of memory for the last segment), the near-peak and combination
grid (about $11$ minutes per half), and the Arb run of piece (F) (about $1.5$ hours on two cores). All other certificates
run in seconds to a few minutes.

\subsection{Trust base}
The proofs rely on:
\begin{enumerate}
\item the written arguments of Sections~\ref{sec:framework}--\ref{sec:third};
\item the standard facts listed in Appendix~\ref{app:facts};
\item the correctness of Arb (every ball returned contains the true value) and of FLINT's exact integer polynomial
arithmetic, as accessed through \texttt{python-flint};
\item the correctness of the Python interpreter executing the certificate programs as written.
\end{enumerate}
IEEE-754 floating point enters only in the superseded float-interval version of piece (F), and \texttt{mpmath} enters only
in consistency checks and in the independent recomputation of the constants of Theorem~\ref{thm:G5} (which were computed
at $60$ digits and independently at $50$ digits).


\appendix
%=====================================================================
\section{Verification and audit record}\label{app:verification}
%=====================================================================

This appendix records, in concise form, how the proof of Theorem~B was checked independently, which defects were found,
and which earlier components were withdrawn. The full reports are in \texttt{third/verification/} of the repository
(\texttt{REFEREE\_SEC3.md}, \texttt{REFEREE\_FULL.md}, \texttt{AUDIT*.md}, \texttt{AUDIT2\_*.md}, \texttt{GAP\_*.md},
\texttt{ARB\_F\_RUN.md}). The proof of Theorem~A is documented by its certificates and consistency checks
(\S\ref{sec:computations}).

\subsection{Independent checks}

\emph{First round.}
\begin{itemize}
\item \emph{\S\ref{sec:G5}} (\texttt{REFEREE\_SEC3.md}): independent re-derivation; own \texttt{mpmath} code at $50$ digits;
$10$ cusps $\times$ $2$ points for~\eqref{eq:G5-grow}/\eqref{eq:G5-ng}; quadrature for Lemma~\ref{lem:bessel-contour};
exhaustive test $\nu\le5000$ plus $1509$ values up to $200000$. Verdict: correct; three presentational gaps (the passage
from the Hankel loop to the vertical line, the arc of $0/1$, the regions of deformation), all filled in \S\ref{sec:G5}.
\item \emph{Band and Laplace, classes $0$--$2$, and coverage} (\texttt{GAP\_AUDIT012.md}): re-derivation, NaN-guarded
rerun of all tiles, $45$ independent exact points. Sound; two latent flaws fixed.
\item \emph{Band, classes $3$--$4$}: an earlier audit's findings are repaired in \texttt{GAP\_BAND.md}.
\item \emph{Laplace, classes $3$--$4$} (\texttt{AUDIT\_LAP34.md}): re-derivation, $6000$-point test of $\varepsilon_g$,
NaN-guarded rerun, $36$ exact points with independent code. Sound after two fixes.
\item \emph{Near peak and combination} (\texttt{AUDIT\_CENTRE.md}): re-derivation of the representation, Euler--Maclaurin
constants, uniformity in $n$, saddle and exact-run logs; exact $P_{5,300}$ comparison. Sound after replacing the
combination step.
\item \emph{Off-peak lemma} (\texttt{AUDIT\_OFFPEAK.md}): coverage in $\theta$, uniformity in $n$,
Lemma~\ref{lem:monoL}, chunk lemma and maximum principle, $312$ exact spot points; Arb version of (F) on $7$ rectangles.
Sound after cosmetic fixes.
\end{itemize}

\emph{Second round} (each auditor wrote their findings before reading the first-round report).
\begin{itemize}
\item \emph{Band, classes $3$--$4$} (\texttt{AUDIT\_BAND.md}): exact run re-run; $16$ values of $n$ recomputed by a direct
product avoiding $E_n$ and the coefficient table; pair and triple counts brute-forced; L1 and $\Psi$ recomputed in
\texttt{mpmath} ($\Psi(10001)=0.19714$). Two cosmetic misstatements. Sound.
\item \emph{Laplace, classes $3$--$4$} (\texttt{AUDIT2\_LAP34.md}): re-derivation from scratch; complete term inventory
including the constants of $T_\rho$ against Theorem~\ref{thm:G5}; monotonicity in $X$ at $6$ levels; a third exact method
(Euler identities modulo primes, CRT), $31$ points up to $n=5000$. No hole. Sound.
\item \emph{Centre} (\texttt{AUDIT2\_CENTRE.md}): Cauchy decomposition and $\cN$ re-derived; $\mu(2.1)=0.103772$; the thin
box recomputed in independent Arb code; monotonicity in $n$ up to $10^5$; quadrature validated to $4\cdot10^{-11}$, then
$120$ points. Findings: the comparison slack $>$ OFF was printed but not asserted (now asserted, Lemma~\ref{lem:comb}); the
binding box $[2.09,2.10]$ is not needed. Sound.
\item \emph{Off-peak} (\texttt{AUDIT2\_OFFPEAK.md}): all ingredients re-derived; (N) compared with the exact $j$-series at
$70$ points; certificates re-run at $N_0=2000$, $10^4$, $10^5$; every (F) rectangle matched; $11$ pre-fix (F) rectangles
re-run; $5570$ adversarial Arb points. Findings: float slivers at $|y|=1.1$ and $L=0.1$; the constant $6.3$ used as a float,
$1.8\cdot10^{-16}$ low. All fixed (\S\ref{sec:offpeak}). Sound.
\item \emph{Classes $0$--$2$} (\texttt{AUDIT2\_CLASS012.md}): representation re-derived, every term inventoried,
monotonicity in $X$ tested up to $1000X_{\min}$, $76$ hard points validated by an exact CRT method over $82$ primes: all
signs correct, true error $15$--$260\times$ below the certified bound. No mathematical error.
\end{itemize}

\emph{Outside referee} (\texttt{REFEREE\_FULL.md}). The whole proof was read end to end, with the certificates treated as
black boxes, $13$ logs and two scripts spot-checked. Verdict: correct with gaps to fill. One major point: for Laplace
classes $0$--$2$ the float-stepped tiling ended at $11.999999999999831$ while the tail cell started at $12$. This is now
closed twice (last tile snapped to $12$, tail cell run from $11.99$; Lemmas~\ref{lem:lap012-tiles}
and~\ref{lem:lap012-tail}). Nine minor points were presentational and are addressed.

\subsection{Defects found and fixed in the components used}
\begin{itemize}
\item \emph{Expansion.} The Farey order must be $N=K(\nu)$, not $K(\nu)-1$; $\Econ$ is rounded up to $27.4636$ (it had been
rounded down); $K_5$ is rounded up to $2.8549955$.
\item \emph{Band.} An earlier majorant used a Bessel envelope in the wrong direction; an earlier extension certificate
undercounted pairs by $4\times$ and triples by $10\times$; an earlier monotonicity script dropped NaN results silently. All
are superseded by the exact run to $n=2000$ and by $\Psi$; the threshold for ``$\Psi$ decreasing'' is $N\ge257$.
\item \emph{Laplace, classes $0$--$2$.} The pass test now rejects NaN (\texttt{not all(x > 0)}); the decay assertion uses
$\kappa-a_5/4$ (the $k=10$ Rankin term) instead of $\kappa-a_5/5$; the $V=12$ sliver is closed.
\item \emph{Laplace, classes $3$--$4$.} The $V=12$ sliver is closed in the same way; NaN pass-through in suprema and in the
pass test is guarded; the log is regenerated with the fixed script.
\item \emph{Centre.} The scalar comparison $|\mathrm{OFF}|\le0.0187\le$ slack of an earlier version is false
($0.018791$ against $0.018703$); it is replaced by the per-box comparison, now asserted in Arb.
\item \emph{Off-peak lemma.} The sliver $(\mathrm{fl}(\pi),\pi]$, the $h$/sign reduction for (N), an inexact key in the
assembly and the validation of the float checker are fixed or superseded by the Arb run of (F); an early exit in the Arb
cell test (which could only reject) was removed; float slivers at $|y|=1.1$ and $L=0.1$ are closed; $6.3$ is an exact ball
everywhere.
\item \emph{Earlier centre certificate.} The strip next to each peak was bounded by nothing; the tail was normalised with a
lower curvature, $7.4\times$ too small; the tail lemma was unproved (now Lemma~\ref{lem:offpeak}); an Euler--Maclaurin
constant was invalid; a Taylor constant was evaluated at the centre instead of over the window; an arctan law was used
without its error; $L<0.05$ was not covered. All are fixed in the present near-peak analysis.
\end{itemize}

\subsection{Withdrawn components}\label{app:withdrawn}
The following earlier components are not used anywhere in this paper.
\begin{itemize}
\item \texttt{B\_uniform5.py}, a $983$-tile Laplace certificate for all classes: NaN cells had been replaced by
$e^{-10^6}$, silently dropping $7168$ cells; with a NaN guard, $63$ of $73$ tiles fail.
\item The old Laplace certificate for classes $3$--$4$ (\texttt{mono\_cert.py}, \texttt{zone2\_cert.py},
\texttt{tail3\_cell.py} and related): it was centred on the kernel-only saddle $m=a_5/W^2$ and treated the $O(N)$ thin
factor $e^{-5nu}$ and the twisted-product factor as amplitudes. Against the exact line integral its main term was off by
factors down to $e^{-13}$, while it certified $B\le0.06$, so ``$B<1$'' proved nothing. Superseded by the re-centred saddle
of \S\ref{sec:lap34}.
\item A folded identity with $D_0=1.12601$: correct, but unnecessary, since \S\ref{sec:t-repr} is exact without folding.
\item The old centre certificate: it had the holes listed above; class $4$ fails at the handover once the tail normalisation
is corrected. Superseded by \S\ref{sec:t-centre}.
\item Earlier band extensions and a $t_{\min}$ table with a wrong normalisation.
\end{itemize}

%=====================================================================
\section{Standard facts and elementary lemmas}\label{app:facts}
%=====================================================================

\subsection{Standard facts}
\begin{fact}[Dedekind eta / partition-function transformation]\label{fact:eta}
This is Fact~\ref{fact:G5-eta}, in Rademacher's normalisation with $hH\equiv-1\pmod k$ \cite[Thm~5.1]{Apostol}. It was
confirmed numerically to $10^{-49}$ at $10$ cusps. The form~\eqref{eq:Gp-transform} for $G_p^\delta$ follows from it
\cite{SZ,Padhi_SZ}.
\end{fact}

\begin{fact}[Farey dissection]\label{fact:farey}
The Farey dissection and the Rademacher path via Ford circles (Lemma~\ref{lem:farey}); see \cite[Ch.~5]{Apostol} and
\cite[Ch.~14]{Rademacher}.
% TODO: verify section numbers (the source cites Apostol Ch. 5, sections 5.4--5.7); Ch. 5 (pp. 94--112 of the 2nd ed.) confirmed,
% section numbers and Thm 5.1 not verified online.
\end{fact}

\begin{fact}[Farey neighbours]\label{fact:neighbours}
Consecutive fractions $h/k$, $h'/k'$ of order $N$ satisfy $k+k'\ge N+1$ \cite[Thm~30]{HardyWright}.
\end{fact}

\begin{fact}[Bessel contour integral]\label{fact:bessel-int}
The Hankel-loop representation of $I_1$ \cite[10.9.19]{DLMF}, \cite[\S6.22]{Watson}. The passage from the loop to the
vertical line is proved in \S\ref{sec:G5}.
\end{fact}

\begin{fact}[Bessel facts]\label{fact:bessel}
$I_\nu(y)>0$ for $y>0$, and $I_\nu(y)$ is decreasing in $\nu\ge0$ for fixed $y>0$ \cite{Cochran,Watson}; the closed forms of
$I_{1/2}$, $I_{3/2}$ and $I_{5/2}$; and $\cP_p(i)=\sqrt{a_p/x}\,I_1(2\sqrt{a_px})$, $x=i-B/2$, is the inverse Laplace
transform of $e^{a_p/u}-1$, which follows from the series of $I_1$. These are used in \S\ref{sec:c-finite},
\S\ref{sec:t-band} and \S\ref{sec:t-repr}.
\end{fact}

\begin{fact}[identities]\label{fact:identities}
$\sum_{c=0}^4\Li_2(\omega_5^cx)=\Li_2(x^5)/5$, with the principal branch for $|x|<1$; Euler's pentagonal theorem; Jacobi's
identity for $(q;q)_\infty^3$.
\end{fact}

\begin{fact}[complex analysis]\label{fact:complex}
Cauchy's theorem, Fubini/Tonelli, the maximum principle for subharmonic functions, and Hadamard's three-circles theorem (the
last two in piece (F)), Poisson summation (Lemma~\ref{lem:c-poisson}).
\end{fact}

\begin{fact}[arithmetic software]\label{fact:software}
Arb, through \texttt{python-flint}: every ball returned contains the true value \cite{Arb}. FLINT \texttt{fmpz\_poly}:
exact integer arithmetic \cite{FLINT}. IEEE-754 correct rounding of $+,-,\times,\div$ is relied on only in the superseded
float-interval version of piece (F).
\end{fact}

\subsection{Elementary lemmas}\label{app:elementary}

\begin{lemma}[Cauchy]\label{lem:cauchy}
If $f$ is holomorphic on $|\zeta-c|\le\rho$ and $|f-f(c)|\le M$ on $|\zeta-c|=\rho$, then
$|f^{(k)}(z)|\le k!M\rho/(\rho-d)^{k+1}$ for $|z-c|\le d<\rho$, $k\ge1$.
\end{lemma}
\begin{proof}
$f^{(k)}(z)=\frac{k!}{2\pi i}\oint(f(\zeta)-f(c))(\zeta-z)^{-k-1}d\zeta$ and $|\zeta-z|\ge\rho-d$.
\end{proof}

\begin{lemma}[midpoint rule]\label{lem:midpoint}
For $g\in C^2[-\frac12,\frac12]$, $|g(0)-\int g|\le\frac18\int|g''|$.
\end{lemma}
\begin{proof}
$g(0)-\int g=-\int K g''$ with $K(t)=\frac12(\frac12-|t|)^2\le\frac18$.
\end{proof}

\begin{lemma}[Koksma, one node per cell]\label{lem:koksma}
If $x_i\in[(i-1)h,ih]$ for $i=1,\dots,M$, then $|h\sum f(x_i)-\int_0^1f|\le h\TV(f)+(1-Mh)\sup|f|$.
\end{lemma}
\begin{proof}
$|hf(x_i)-\int_{C_i}f|\le h\TV_{C_i}(f)$; sum, and bound the uncovered part.
\end{proof}


\bibliographystyle{amsplain}
\bibliography{refs}

\end{document}
